\chapter{Многомерный анализ}
\label{chap:multivariable-calculus}

Многомерный анализ изучает отображения, у которых входом или выходом служит сразу несколько чисел. Вектор параметров может описывать геометрию детали, состояние физической системы, настройки алгоритма или признаки объекта. Значение отображения может быть одной характеристикой либо новым вектором измерений. Поэтому прежде чем вводить производные, необходимо точно определить, что означает близость точек, предел последовательности, непрерывность модели и замкнутость допустимой области.

В этой главе теория излагается с опорой на курс В.~А.~Зорича. Из общего топологического аппарата сохраняется только то, что непосредственно требуется для анализа отображений в конечномерных пространствах, оптимизации и последующих вычислений.

\section{Пространство \texorpdfstring{$\R^n$}{R n}, предел и непрерывность}
\label{sec:rn-limit-continuity}

\subsection{Точки, векторы, норма и расстояние}

Пространство
\[
\R^n
=
\{x=(x_1,\dots,x_n)\mid x_i\in\R\}
\]
можно рассматривать одновременно как множество точек и как векторное пространство. В геометрической задаче элемент \(x\) обозначает положение или набор параметров. Разность \(y-x\) является вектором изменения, переводящим точку \(x\) в точку \(y\).

\begin{definition}
Евклидовой нормой вектора \(x=(x_1,\dots,x_n)\in\R^n\) называется число
\[
\norm{x}
=
\sqrt{x_1^2+\cdots+x_n^2}.
\]
Евклидово расстояние между точками \(x,y\in\R^n\) определяется формулой
\[
d(x,y)=\norm{x-y}.
\]
\end{definition}

Норма измеряет величину изменения, а расстояние --- различие двух состояний. Например, если
\[
x=(T,p,v)
\]
содержит нормированные температуру, давление и скорость, то \(\norm{x-y}\) объединяет изменения трёх компонент в одну величину. Нормирование существенно: складывать квадрат метра с квадратом паскаля без предварительного выбора масштаба физически бессмысленно.

Евклидова норма обладает свойствами
\[
\norm{x}\geq0,
\qquad
\norm{x}=0\iff x=0,
\qquad
\norm{\lambda x}=\abs{\lambda}\norm{x},
\]
\[
\norm{x+y}\leq\norm{x}+\norm{y}.
\]
Последнее соотношение называется неравенством треугольника. Из него следует полезная оценка
\[
\abs{\norm{x}-\norm{y}}
\leq
\norm{x-y}.
\]
Она показывает, что малое изменение самого вектора не может вызвать большое изменение его длины.

\begin{remark}
На \(\R^n\) можно задавать и другие нормы, например
\[
\norm{x}_1=\sum_{i=1}^n\abs{x_i},
\qquad
\norm{x}_{\infty}=\max_{1\leq i\leq n}\abs{x_i}.
\]
В конечномерном пространстве они приводят к одному и тому же понятию сходимости и непрерывности. Далее, если не указано иное, используется евклидова норма.
\end{remark}

\subsection{Шары и окрестности}

\begin{definition}
Открытым шаром с центром \(a\in\R^n\) и радиусом \(r>0\) называется множество
\[
B_r(a)
=
\{x\in\R^n\mid \norm{x-a}<r\}.
\]
Замкнутым шаром называется множество
\[
\overline B_r(a)
=
\{x\in\R^n\mid \norm{x-a}\leq r\}.
\]
\end{definition}

Условие \(x\in B_r(a)\) означает, что отклонение \(x-a\) меньше заданного допуска \(r\). Поэтому шар можно читать как множество всех возмущений рабочей точки, допустимых по выбранной норме.

\begin{definition}
Окрестностью точки \(a\in\R^n\) называется любое открытое множество, содержащее \(a\). В частности, каждый шар \(B_r(a)\) является окрестностью точки \(a\).
\end{definition}

В определениях предела достаточно пользоваться шарами. Формулировка через произвольные окрестности понадобится лишь для геометрической интерпретации.

\subsection{Открытые и замкнутые множества}

\begin{definition}
Множество \(G\subseteq\R^n\) называется открытым, если для каждой точки \(a\in G\) существует радиус \(r>0\), такой что
\[
B_r(a)\subseteq G.
\]
\end{definition}

Открытость означает наличие у каждой допустимой точки некоторого запаса: достаточно малое возмущение не выводит её из множества. Радиус запаса может зависеть от точки.

\begin{example}
Множество
\[
G=
\{(x_1,x_2)\in\R^2\mid x_1^2+x_2^2<1\}
\]
открыто. Если \(a\in G\), то \(\norm{a}<1\). Выберем
\[
r=\frac{1-\norm{a}}{2}>0.
\]
Для любой точки \(x\in B_r(a)\) по неравенству треугольника
\[
\norm{x}
\leq
\norm{x-a}+\norm{a}
<
\frac{1-\norm{a}}{2}+\norm{a}
<1.
\]
Следовательно, \(B_r(a)\subseteq G\).
\end{example}

\begin{definition}
Множество \(F\subseteq\R^n\) называется замкнутым, если его дополнение
\[
\R^n\setminus F
\]
открыто.
\end{definition}

Замкнутое множество содержит точки, получающиеся как пределы допустимых последовательностей. Это свойство будет сформулировано точно после введения сходимости.

\begin{claim}
Произвольное объединение открытых множеств открыто, а пересечение конечного числа открытых множеств открыто. Произвольное пересечение замкнутых множеств замкнуто, а объединение конечного числа замкнутых множеств замкнуто.
\end{claim}

Ограничение на конечность существенно. Например,
\[
G_k=\left(-\frac1k,\frac1k\right)
\]
открыто в \(\R\) для каждого \(k\), но
\[
\bigcap_{k=1}^{\infty}G_k=\{0\}
\]
не является открытым множеством.

\subsection{Внутренние, граничные и предельные точки}

\begin{definition}
Пусть \(E\subseteq\R^n\). Точка \(a\in\R^n\) называется:

Точку \(a\) называют внутренней для \(E\), если некоторый шар \(B_r(a)\) целиком лежит в \(E\); внешней, если некоторый такой шар не пересекает \(E\); граничной, если каждый шар с центром в \(a\) содержит точки как из \(E\), так и из \(\R^n\setminus E\).
\end{definition}

Множество внутренних точек обозначается через \(\operatorname{int}E\), а множество граничных точек --- через \(\partial E\). Граница показывает, где исчезает запас допустимости: сколь угодно малое изменение может перевести точку с одной стороны ограничения на другую.

\begin{definition}
Точка \(a\in\R^n\) называется предельной точкой множества \(E\), если для каждого \(r>0\) существует точка
\[
x\in E,
\qquad
x\neq a,
\qquad
\norm{x-a}<r.
\]
\end{definition}

Предельная точка не обязана принадлежать самому множеству. Например, точка \(0\) является предельной для интервала \((0,1)\), но не лежит в нём.

\begin{definition}
Замыканием множества \(E\) называется объединение \(E\) со всеми его предельными точками. Замыкание обозначается через \(\overline E\).
\end{definition}

\begin{criterion}[критерий замкнутости]
Множество \(F\subseteq\R^n\) замкнуто тогда и только тогда, когда оно содержит все свои предельные точки.
\end{criterion}

Для практической области параметров это означает: если допустимые состояния можно приближать допустимыми состояниями сколь угодно точно, то предельное состояние также должно считаться допустимым.

\subsection{Последовательности в \texorpdfstring{$\R^n$}{R n}}

\begin{definition}
Последовательность точек \(x^{(1)},x^{(2)},\dots\in\R^n\) сходится к точке \(a\in\R^n\), если для каждого \(\varepsilon>0\) существует номер \(N\), такой что
\[
k\geq N
\quad\Longrightarrow\quad
\norm{x^{(k)}-a}<\varepsilon.
\]
В этом случае пишут
\[
\lim_{k\to\infty}x^{(k)}=a
\quad\text{или}\quad
x^{(k)}\to a.
\]
\end{definition}

Индекс последовательности записан в скобках, чтобы не путать его с номером координаты. Если
\[
x^{(k)}=
\bigl(x_1^{(k)},\dots,x_n^{(k)}\bigr),
\qquad
a=(a_1,\dots,a_n),
\]
то сходимость можно проверять по координатам.

\begin{theorem}[покоординатный критерий]
Последовательность \(x^{(k)}\in\R^n\) сходится к \(a\in\R^n\) тогда и только тогда, когда
\[
x_i^{(k)}\to a_i
\]
для каждой координаты \(i=1,\dots,n\).
\end{theorem}

\begin{proofidea}
Из
\[
\abs{x_i^{(k)}-a_i}
\leq
\norm{x^{(k)}-a}
\]
следует, что сходимость в норме влечёт покоординатную сходимость. Обратно, если каждая координатная разность стремится к нулю, то при достаточно большом \(k\)
\[
\abs{x_i^{(k)}-a_i}<\frac{\varepsilon}{\sqrt n}
\]
для всех \(i\). Тогда
\[
\norm{x^{(k)}-a}
=
\sqrt{\sum_{i=1}^n\abs{x_i^{(k)}-a_i}^2}
<\varepsilon.
\]
\end{proofidea}

\begin{criterion}[последовательный критерий замкнутости]
Множество \(F\subseteq\R^n\) замкнуто тогда и только тогда, когда из условий
\[
x^{(k)}\in F,
\qquad
x^{(k)}\to a
\]
следует \(a\in F\).
\end{criterion}

Этот критерий удобен для ограничений, заданных равенствами и нестрогими неравенствами. Если при предельном переходе каждое ограничение сохраняется, область замкнута.

\subsection{Предел отображения}

Пусть \(E\subseteq\R^n\), а отображение
\[
f:E\to\R^m
\]
сопоставляет входному вектору \(x\) выходной вектор \(f(x)\). Точка \(a\) может не принадлежать \(E\), но должна быть предельной точкой области определения.

\begin{definition}
Вектор \(b\in\R^m\) называется пределом отображения \(f\) при \(x\to a\), если для каждого \(\varepsilon>0\) существует \(\delta>0\), такое что из условий
\[
x\in E,
\qquad
0<\norm{x-a}<\delta
\]
следует
\[
\norm{f(x)-b}<\varepsilon.
\]
Используется запись
\[
\lim_{x\to a}f(x)=b.
\]
\end{definition}

Число \(\varepsilon\) задаёт требуемую точность выхода, а \(\delta\) --- достаточную точность входа. Зависимость \(\delta\) от \(\varepsilon\) выражает устойчивость модели вблизи точки \(a\).

\begin{theorem}[критерий Гейне]
Равенство
\[
\lim_{x\to a}f(x)=b
\]
выполнено тогда и только тогда, когда для любой последовательности
\[
x^{(k)}\in E\setminus\{a\},
\qquad
x^{(k)}\to a,
\]
имеет место
\[
f\bigl(x^{(k)}\bigr)\to b.
\]
\end{theorem}

Критерий особенно полезен для опровержения существования предела. Достаточно найти два способа приближения к одной точке, дающие разные пределы значений.

\begin{example}
Рассмотрим функцию
\[
f(x,y)=
\frac{xy}{x^2+y^2},
\qquad
(x,y)\neq(0,0).
\]
По прямой \(y=x\) получаем
\[
f(x,x)=\frac12,
\]
а по прямой \(y=0\)
\[
f(x,0)=0.
\]
Следовательно, при \((x,y)\to(0,0)\) единого предела нет. Проверка только одного пути была бы недостаточной.
\end{example}

\subsection{Непрерывность}

\begin{definition}
Отображение \(f:E\to\R^m\) называется непрерывным в точке \(a\in E\), если
\[
\lim_{x\to a}f(x)=f(a).
\]
Эквивалентно: для каждого \(\varepsilon>0\) существует \(\delta>0\), такое что
\[
x\in E,
\qquad
\norm{x-a}<\delta
\quad\Longrightarrow\quad
\norm{f(x)-f(a)}<\varepsilon.
\]
\end{definition}

Отличие от определения предела состоит в том, что точка \(x=a\) теперь допускается, а предельное значение обязано совпадать с фактическим значением модели.

\begin{theorem}[последовательный критерий непрерывности]
Отображение \(f:E\to\R^m\) непрерывно в точке \(a\in E\) тогда и только тогда, когда из
\[
x^{(k)}\in E,
\qquad
x^{(k)}\to a
\]
следует
\[
f\bigl(x^{(k)}\bigr)\to f(a).
\]
\end{theorem}

Если
\[
f=(f_1,\dots,f_m),
\]
то \(f\) непрерывно тогда и только тогда, когда непрерывна каждая координатная функция \(f_j\). Суммы, произведения и композиции непрерывных функций непрерывны там, где соответствующие операции определены. Поэтому многочлены непрерывны на всём \(\R^n\), а рациональные функции --- во всех точках, где знаменатель не обращается в нуль.

\begin{application}[устойчивость вычислительной модели]
Пусть \(x\) --- вектор измеренных параметров, а \(f(x)\) --- вычисляемая характеристика. Непрерывность в рабочей точке \(a\) означает, что для любой допустимой ошибки результата можно подобрать допустимую ошибку входных данных. Она не сообщает скорость роста ошибки; количественные оценки появятся после введения дифференциала и якобиана.
\end{application}

\subsection{Ограниченность и компактность}

\begin{definition}
Множество \(E\subseteq\R^n\) называется ограниченным, если существуют точка \(a\in\R^n\) и число \(R>0\), для которых
\[
E\subseteq B_R(a).
\]
\end{definition}

Ограниченность означает, что параметры не могут уходить сколь угодно далеко. Замкнутость и ограниченность решают разные задачи: первое свойство сохраняет предельные состояния, второе исключает неограниченный уход.

\begin{definition}
Множество \(K\subseteq\R^n\) называется компактным, если из любой последовательности точек \(x^{(k)}\in K\) можно выбрать подпоследовательность, сходящуюся к некоторой точке \(a\in K\).
\end{definition}

\begin{theorem}[Гейне--Бореля]
Множество \(K\subseteq\R^n\) компактно тогда и только тогда, когда оно замкнуто и ограничено.
\end{theorem}

В конечномерной задаче этот критерий обычно проще самого определения. Например, область, заданная конечным набором непрерывных ограничений
\[
g_i(x)\leq0,
\qquad
h_j(x)=0,
\]
замкнута, а явные верхние и нижние границы параметров дают ограниченность.

\begin{theorem}[Вейерштрасс]
Если \(K\subseteq\R^n\) компактно, а функция \(f:K\to\R\) непрерывна, то существуют точки \(x_{\min},x_{\max}\in K\), для которых
\[
f(x_{\min})
=
\min_{x\in K}f(x),
\qquad
f(x_{\max})
=
\max_{x\in K}f(x).
\]
\end{theorem}

Теорема гарантирует существование оптимального решения, но сама по себе не даёт алгоритма его поиска. В последующих разделах градиент, гессиан и условия на границе будут использоваться для нахождения кандидатов.

\begin{problem}
\textit{Задача по материалу Зорича.}
Рассмотрим множество
\[
D=
\left\{(x,y)\in\R^2
\;\middle|\;
 x^2+y^2\leq1,
\ y\geq0
\right\}.
\]
Найти его внутренность и границу, доказать замкнутость и компактность.
\end{problem}

\begin{proof}[Решение]
Множество \(D\) является верхним замкнутым полукругом вместе с диаметром.

Точка лежит во внутренности тогда и только тогда, когда оба ограничения выполняются строго:
\[
\operatorname{int}D
=
\left\{(x,y)\in\R^2
\;\middle|\;
 x^2+y^2<1,
\ y>0
\right\}.
\]
Действительно, при строгих неравенствах остаётся положительный запас до окружности и до прямой \(y=0\), поэтому достаточно малый шар остаётся внутри \(D\). Если хотя бы одно ограничение обращается в равенство, любой шар содержит точки, нарушающие это ограничение.

Следовательно, граница состоит из верхней полуокружности и диаметра:
\[
\partial D
=
\left\{(x,y)\mid x^2+y^2=1,\ y\geq0\right\}
\cup
\left\{(x,0)\mid -1\leq x\leq1\right\}.
\]

Покажем замкнутость последовательным критерием. Пусть
\[
(x_k,y_k)\in D,
\qquad
(x_k,y_k)\to(x,y).
\]
Тогда покоординатно \(x_k\to x\), \(y_k\to y\). Непрерывность многочлена даёт
\[
x_k^2+y_k^2\to x^2+y^2.
\]
Из неравенств \(x_k^2+y_k^2\leq1\) и \(y_k\geq0\) при предельном переходе получаем
\[
x^2+y^2\leq1,
\qquad
y\geq0.
\]
Значит, \((x,y)\in D\), и \(D\) замкнуто.

Кроме того,
\[
\norm{(x,y)}\leq1
\]
для всех \((x,y)\in D\), поэтому множество ограничено. По теореме Гейне--Бореля \(D\) компактно.
\end{proof}

\begin{problem}
\textit{Практическая задача: допустимая область проектирования.}
Параметры прямоугольной панели --- ширина \(w\) и высота \(h\), измеряемые в метрах. Технологические ограничения имеют вид
\[
1\leq w\leq3,
\qquad
1\leq h\leq2,
\qquad
w+h\leq4.
\]
Качество выбранных размеров описывается непрерывной функцией
\[
Q(w,h)
=
10-(w-2)^2-2(h-1.5)^2.
\]
Доказать существование наилучшего варианта и найти его. Объяснить смысл каждой части модели.
\end{problem}

\begin{proof}[Решение]
Входом модели является вектор параметров
\[
x=(w,h)\in\R^2.
\]
Первые два двойных неравенства задают технологические диапазоны ширины и высоты. Ограничение \(w+h\leq4\) исключает слишком крупные сочетания размеров. Допустимое множество равно
\[
K=
\left\{(w,h)\in\R^2
\;\middle|\;
1\leq w\leq3,
\ 1\leq h\leq2,
\ w+h\leq4
\right\}.
\]

Каждое ограничение нестрогое и задаётся непрерывной функцией, поэтому \(K\) замкнуто. Из диапазонов \(1\leq w\leq3\), \(1\leq h\leq2\) следует ограниченность. Следовательно, \(K\) компактно.

Функция \(Q\) является многочленом и потому непрерывна. Теорема Вейерштрасса гарантирует, что максимум на \(K\) достигается.

Вычитаемые слагаемые являются неотрицательными квадратами, поэтому
\[
Q(w,h)\leq10.
\]
Равенство возможно при
\[
w=2,
\qquad
h=1.5.
\]
Проверяем ограничения:
\[
1\leq2\leq3,
\qquad
1\leq1.5\leq2,
\qquad
2+1.5=3.5\leq4.
\]
Рабочая точка допустима, поэтому
\[
Q_{\max}=10
\quad\text{при}\quad
(w,h)=(2,1.5).
\]
Квадраты штрафуют отклонения от размеров \(2\) м и \(1.5\) м, причём второе отклонение имеет вдвое больший вес. Теорема Вейерштрасса обосновывает существование решения, а вид \(Q\) позволяет найти его без дифференцирования.
\end{proof}

% =============================================================================
% V02 — differentiability, differential, Jacobian and gradient
% =============================================================================

\section{Дифференцируемость, дифференциал, якобиан и градиент}
\label{sec:differential-jacobian-gradient}

Непрерывность отвечает на качественный вопрос: остаётся ли выход модели близким при малом изменении входа. Дифференцируемость даёт более сильный результат --- она выделяет линейную часть этого изменения. Именно эта линейная часть позволяет предсказывать чувствительность, строить касательные объекты, оценивать влияние параметров и далее применять правило цепочки.

Пусть
\[
F\colon G\subseteq\R^n\to\R^m,
\qquad
x=(x_1,\dots,x_n),
\qquad
F(x)=\bigl(F_1(x),\dots,F_m(x)\bigr),
\]
где \(G\) --- открытое множество. Вектор \(x\) содержит \(n\) входных переменных, а \(F(x)\) --- \(m\) выходных величин. Для рабочей точки \(a\in G\) и малого приращения \(h\in\R^n\) рассматривается изменение
\[
\Delta F
=
F(a+h)-F(a).
\]
Задача дифференциального исчисления состоит в том, чтобы отделить главную, линейную по \(h\), часть этого изменения от остатка более высокого порядка малости.

\subsection{Дифференцируемость как локальная линейная модель}

\begin{definition}
Отображение \(F\colon G\subseteq\R^n\to\R^m\) называется дифференцируемым в точке \(a\in G\), если существует линейное отображение
\[
L\colon\R^n\to\R^m
\]
такое, что при \(h\to0\)
\begin{equation}
F(a+h)
=
F(a)+Lh+r(h),
\qquad
\frac{\norm{r(h)}}{\norm{h}}\to0.
\label{eq:differentiability-definition}
\end{equation}
Линейное отображение \(L\) называется дифференциалом, или производной, отображения \(F\) в точке \(a\) и обозначается
\[
DF(a).
\]
Его матрица в стандартных базисах обозначается через \(\Jac{F}(a)\).
\end{definition}

Условие
\[
\frac{\norm{r(h)}}{\norm{h}}\to0
\]
записывают также как
\[
r(h)=o(\norm{h}).
\]
Оно означает, что остаток мал не просто сам по себе, а мал по сравнению с размером входного возмущения. Поэтому формула
\begin{equation}
F(a+h)-F(a)
\approx
DF(a)h
\label{eq:linear-response}
\end{equation}
становится тем точнее, чем меньше \(\norm{h}\).

В формуле \eqref{eq:differentiability-definition} роли объектов различны. Точка \(a\) задаёт рабочий режим. Вектор \(h\) задаёт конкретное изменение входа. Линейный оператор \(DF(a)\) зависит от рабочей точки, но после её фиксации действует на все малые приращения \(h\). Вектор \(DF(a)h\) предсказывает первое изменение выхода, а \(r(h)\) содержит нелинейную поправку.

\begin{theorem}[единственность дифференциала]
Если отображение \(F\) дифференцируемо в точке \(a\), то линейное отображение \(DF(a)\) определяется единственным образом.
\end{theorem}

\begin{proofidea}
Предположим, что линейную главную часть задают два оператора \(L_1\) и \(L_2\). Для фиксированного вектора \(v\) положим \(h=t v\). Вычитая два представления приращения и деля на \(t\), получаем
\[
(L_1-L_2)v\to0
\qquad
\text{при }t\to0.
\]
Левая часть от \(t\) не зависит, поэтому \((L_1-L_2)v=0\). Так как это верно для любого \(v\), операторы совпадают.
\end{proofidea}

\begin{theorem}
Если отображение \(F\) дифференцируемо в точке \(a\), то оно непрерывно в этой точке.
\end{theorem}

\begin{proof}
Из определения дифференцируемости
\[
F(a+h)-F(a)=DF(a)h+r(h).
\]
Линейное отображение в конечномерных пространствах непрерывно, поэтому существует константа \(C\geq0\), для которой
\[
\norm{DF(a)h}\leq C\norm{h}.
\]
Следовательно,
\[
\norm{F(a+h)-F(a)}
\leq
C\norm{h}+\norm{r(h)}.
\]
Оба слагаемых стремятся к нулю при \(h\to0\), значит,
\[
F(a+h)\to F(a).
\]
\end{proof}

Обратное утверждение неверно: непрерывность не гарантирует существования единой линейной модели. Например, функция \(f(x)=\abs{x}\) непрерывна при \(x=0\), но её левый и правый линейные наклоны различны.

\subsection{Частные производные: изменение одной входной переменной}

Обозначим через \(e_j\) \(j\)-й стандартный базисный вектор пространства \(\R^n\):
\[
e_1=(1,0,\dots,0),
\quad\dots,\quad
e_n=(0,\dots,0,1).
\]
Приращение \(t e_j\) изменяет только координату \(x_j\), оставляя остальные входы фиксированными.

\begin{definition}
Частной производной отображения \(F\) по переменной \(x_j\) в точке \(a\) называется предел
\begin{equation}
\pdv{F}{x_j}(a)
=
\lim_{t\to0}
\frac{F(a+t e_j)-F(a)}{t},
\label{eq:partial-vector-map}
\end{equation}
если он существует.
\end{definition}

Если \(F\) принимает значения в \(\R^m\), то частная производная \(\partial F/\partial x_j(a)\) является вектором из \(\R^m\). Его \(i\)-я компонента равна
\[
\pdv{F_i}{x_j}(a).
\]
Таким образом, частная производная отвечает на вопрос: как одновременно меняются все выходы, когда изменяется только один конкретный вход.

\begin{theorem}
Если \(F\) дифференцируемо в точке \(a\), то в этой точке существуют все частные производные и
\begin{equation}
\pdv{F}{x_j}(a)=DF(a)e_j.
\label{eq:partial-as-column}
\end{equation}
Для произвольного приращения
\[
h=h_1e_1+\cdots+h_ne_n
\]
выполняется
\begin{equation}
DF(a)h
=
\sum_{j=1}^{n}
\pdv{F}{x_j}(a)h_j.
\label{eq:differential-from-partials}
\end{equation}
\end{theorem}

Формула \eqref{eq:partial-as-column} объясняет устройство матрицы Якоби: её \(j\)-й столбец должен описывать реакцию всех выходов на изменение \(j\)-го входа.

Существование всех частных производных ещё не означает дифференцируемости. Это принципиальное отличие от функции одной переменной.

\begin{counterexample}[по В.~А.~Зоричу]
Рассмотрим функцию
\[
f(x,y)
=
\begin{cases}
0, & xy=0,\\
1, & xy\neq0.
\end{cases}
\]
На обеих координатных осях функция равна нулю. Поэтому
\[
\pdv{f}{x}(0,0)
=
\lim_{t\to0}\frac{f(t,0)-f(0,0)}{t}=0,
\]
\[
\pdv{f}{y}(0,0)
=
\lim_{t\to0}\frac{f(0,t)-f(0,0)}{t}=0.
\]
Однако вдоль последовательности \((1/k,1/k)\to(0,0)\) значения функции равны единице. Следовательно, функция не непрерывна в начале координат и потому не может быть там дифференцируема.
\end{counterexample}

На практике дифференцируемость часто устанавливают по следующему достаточному условию.

\begin{theorem}[достаточное условие дифференцируемости]
Пусть все частные производные компонент отображения
\[
F\colon G\subseteq\R^n\to\R^m
\]
существуют в некоторой окрестности точки \(a\) и непрерывны в точке \(a\). Тогда \(F\) дифференцируемо в точке \(a\).
\end{theorem}

Условие является достаточным, но не необходимым: дифференцируемая функция может иметь частные производные, которые не являются непрерывными в окрестности. Для вычислительных моделей, составленных из многочленов, экспонент, логарифмов и тригонометрических функций в области их определения, непрерывность частных производных обычно проверяется непосредственно.

\subsection{Производная по направлению}

Частная производная исследует движение только вдоль координатной оси. В прикладной задаче параметры часто меняются совместно, поэтому нужен произвольный вектор направления.

\begin{definition}
Производной отображения \(F\) в точке \(a\) по вектору \(v\in\R^n\) называется предел
\begin{equation}
D_vF(a)
=
\lim_{t\to0}
\frac{F(a+t v)-F(a)}{t},
\label{eq:directional-derivative}
\end{equation}
если он существует.
\end{definition}

Параметр \(t\) задаёт величину перемещения, а вектор \(v\) --- соотношение, в котором меняются координаты. Если \(\norm{v}=1\), производная показывает изменение выхода на единицу пути в направлении \(v\). Если \(v\) не единичный, его длина дополнительно масштабирует скорость изменения.

\begin{theorem}
Если \(F\) дифференцируемо в точке \(a\), то производная по любому вектору \(v\) существует и
\begin{equation}
D_vF(a)=DF(a)v.
\label{eq:directional-via-differential}
\end{equation}
\end{theorem}

\begin{proof}
Подставим в определение дифференцируемости \(h=t v\):
\[
F(a+t v)-F(a)
=
t\,DF(a)v+r(tv).
\]
После деления на \(t\neq0\) получаем
\[
\frac{F(a+t v)-F(a)}{t}
=
DF(a)v+\frac{r(tv)}{t}.
\]
Так как \(r(tv)=o(\abs{t}\norm{v})\), последнее слагаемое стремится к нулю.
\end{proof}

Формула \eqref{eq:directional-via-differential} показывает, что дифференциал содержит сразу все производные по направлениям. Отдельное вычисление предела вдоль каждой прямой не требуется. При этом существование производных по всем направлениям само по себе не гарантирует дифференцируемости: направленные пределы могут не складываться в одну согласованную линейную модель.

\subsection{Матрица Якоби}

Линейное отображение \(DF(a)\colon\R^n\to\R^m\) после выбора стандартных базисов задаётся матрицей размера \(m\times n\).

\begin{definition}
Матрицей Якоби отображения
\[
F=(F_1,\dots,F_m)\colon G\subseteq\R^n\to\R^m
\]
в точке \(a\) называется матрица
\begin{equation}
\Jac{F}(a)
=
\begin{pmatrix}
\dfrac{\partial F_1}{\partial x_1}(a)
& \cdots &
\dfrac{\partial F_1}{\partial x_n}(a)\\[1.1em]
\vdots & \ddots & \vdots\\[0.4em]
\dfrac{\partial F_m}{\partial x_1}(a)
& \cdots &
\dfrac{\partial F_m}{\partial x_n}(a)
\end{pmatrix}.
\label{eq:jacobian-definition}
\end{equation}
\end{definition}

Строка \(i\) соответствует выходу \(F_i\), столбец \(j\) --- входу \(x_j\), а элемент \((i,j)\) измеряет локальную чувствительность этого выхода к этому входу. Поэтому число строк равно числу выходов, а число столбцов --- числу входов.

Если векторы записываются столбцами, то действие дифференциала имеет вид
\begin{equation}
DF(a)h
=
\Jac{F}(a)h.
\label{eq:differential-jacobian-action}
\end{equation}
Следовательно, линейная аппроксимация записывается как
\begin{equation}
F(a+h)
\approx
F(a)+\Jac{F}(a)h.
\label{eq:jacobian-linearization}
\end{equation}

Матрица Якоби не является новым объектом, независимым от дифференциала. Дифференциал --- линейное отображение, а якобиан --- его матрица в выбранных координатах. При замене базисов матрица меняется, но само линейное отображение остаётся тем же.

\begin{remark}[о термине «якобиан»]
Термином «якобиан» называют и матрицу \(\Jac{F}\), и определитель \(\det \Jac{F}\), когда \(m=n\). В этой главе словосочетание «матрица Якоби» всегда означает матрицу, а «определитель Якоби» --- её определитель.
\end{remark}

\begin{caution}[единицы измерения]
Элементы одной матрицы Якоби могут иметь разные физические единицы. Если \(F_i\) измеряется, например, в килограммах, а \(x_j\) --- в сантиметрах, то
\[
\pdv{F_i}{x_j}
\]
измеряется в килограммах на сантиметр. Поэтому сравнивать численные величины разных столбцов без масштабирования входных параметров не всегда корректно.
\end{caution}

\subsection{Дифференциал скалярной функции и градиент}

Пусть теперь выход один:
\[
f\colon G\subseteq\R^n\to\R.
\]
Матрица Якоби имеет одну строку,
\[
\Jac{f}(a)
=
\begin{pmatrix}
\pdv{f}{x_1}(a)&\cdots&\pdv{f}{x_n}(a)
\end{pmatrix}.
\]
Те же коэффициенты удобно собрать в столбец.

\begin{definition}
Градиентом дифференцируемой функции \(f\colon G\subseteq\R^n\to\R\) в точке \(a\) называется вектор
\begin{equation}
\grad f(a)
=
\begin{pmatrix}
\dfrac{\partial f}{\partial x_1}(a)\\[0.8em]
\vdots\\[0.3em]
\dfrac{\partial f}{\partial x_n}(a)
\end{pmatrix}.
\label{eq:gradient-definition}
\end{equation}
\end{definition}

Дифференциал является линейным функционалом, поэтому
\begin{equation}
Df(a)h
=
\inner{\grad f(a)}{h}
=
\grad f(a)^{\trans}h.
\label{eq:differential-gradient}
\end{equation}
Таким образом, якобиан скалярной функции является строкой, а градиент --- транспонированным к ней столбцом:
\[
\Jac{f}(a)=\grad f(a)^{\trans}.
\]

Для единичного направления \(v\) получаем
\begin{equation}
D_vf(a)
=
\inner{\grad f(a)}{v}.
\label{eq:directional-gradient}
\end{equation}
По неравенству Коши--Буняковского
\[
D_vf(a)
\leq
\norm{\grad f(a)}\norm{v}
=
\norm{\grad f(a)}.
\]
Если \(\grad f(a)\neq0\), наибольшее значение достигается при
\[
v=\frac{\grad f(a)}{\norm{\grad f(a)}}.
\]
Следовательно, градиент указывает направление наиболее быстрого локального роста функции, а \(-\grad f(a)\) --- направление наиболее быстрого убывания среди единичных направлений.

Геометрический смысл проявляется на поверхности уровня
\[
f(x)=c.
\]
Если гладкая кривая \(\gamma(t)\) лежит на этой поверхности и проходит через точку \(a=\gamma(0)\), то функция \(f(\gamma(t))\) постоянна. Дифференцируя по \(t\), получаем
\[
\inner{\grad f(a)}{\gamma'(0)}=0.
\]
Значит, градиент ортогонален касательным направлениям к поверхности уровня. В пространстве \(\R^3\) касательная плоскость к уровню \(f(x,y,z)=c\) в точке \(a\) задаётся уравнением
\[
\inner{\grad f(a)}{x-a}=0,
\]
если \(\grad f(a)\neq0\).

При построении локальной модели сначала задают смысл входов \(x\in\R^n\) и выходов \(F(x)\in\R^m\). Затем вычисляют производные для каждой пары «выход--вход», размещают их в матрице \(\Jac{F}(a)\), подставляют рабочую точку и находят \(\Jac{F}(a)h\). Полученное приращение проверяют по размерности и, когда возможно, сравнивают с точным пересчётом модели.

\begin{problem}[Задача по материалу В.~А.~Зорича: скалярная функция]
Для функции
\[
f(x,y)=x^2e^y+xy^2
\]
в точке
\[
a=(1,0)
\]
найти частные производные, дифференциал, градиент и производную по единичному направлению
\[
v=\left(\frac35,\frac45\right).
\]
Записать линейную аппроксимацию и проверить производную по направлению непосредственным сведением к функции одной переменной.
\end{problem}

\begin{proof}[Решение]
Входом является точка \((x,y)\in\R^2\), выходом --- одно число \(f(x,y)\). Поэтому дифференциал действует из \(\R^2\) в \(\R\), матрица Якоби имеет размер \(1\times2\), а градиент имеет две компоненты.

Сначала вычисляем производную по первому входу \(x\), считая \(y\) постоянным:
\[
\pdv{f}{x}(x,y)
=
2xe^y+y^2.
\]
Производная по второму входу \(y\), при фиксированном \(x\), равна
\[
\pdv{f}{y}(x,y)
=
x^2e^y+2xy.
\]
Подставляем рабочую точку \(a=(1,0)\):
\[
\pdv{f}{x}(1,0)
=
2\cdot1\cdot e^0+0^2
=2,
\]
\[
\pdv{f}{y}(1,0)
=
1^2e^0+2\cdot1\cdot0
=1.
\]
Следовательно,
\[
\Jac{f}(1,0)
=
\begin{pmatrix}2&1\end{pmatrix},
\qquad
\grad f(1,0)
=
\begin{pmatrix}2\\1\end{pmatrix}.
\]
Для приращения \(h=(h_1,h_2)\) дифференциал равен
\[
Df(1,0)h
=
\begin{pmatrix}2&1\end{pmatrix}
\begin{pmatrix}h_1\\h_2\end{pmatrix}
=
2h_1+h_2.
\]
Поскольку
\[
f(1,0)=1,
\]
локальная линейная модель имеет вид
\begin{equation}
 f(1+h_1,h_2)
=
1+2h_1+h_2+o\!\left(\sqrt{h_1^2+h_2^2}\right).
\label{eq:scalar-task-linearization}
\end{equation}

Проверяем, что направление единичное:
\[
\norm{v}
=
\sqrt{\left(\frac35\right)^2+\left(\frac45\right)^2}
=1.
\]
По формуле через градиент
\[
D_vf(1,0)
=
\inner{\grad f(1,0)}{v}
=
2\cdot\frac35+1\cdot\frac45
=2.
\]

Проверим результат непосредственно. Ограничим функцию на прямую
\[
(x,y)
=
(1,0)+t\left(\frac35,\frac45\right)
=
\left(1+\frac{3t}{5},\frac{4t}{5}\right).
\]
Получаем функцию одной переменной
\[
\varphi(t)
=
f\left(1+\frac{3t}{5},\frac{4t}{5}\right).
\]
По обычному правилу дифференцирования композиции в точке \(t=0\)
\[
\varphi'(0)
=
\pdv{f}{x}(1,0)\frac35
+
\pdv{f}{y}(1,0)\frac45
=
2.
\]
Это совпадает с вычислением через градиент.
\end{proof}

\begin{problem}[Практическая задача: геометрическая модель]
Прямоугольная алюминиевая пластина описывается шириной \(w\), высотой \(h\) и толщиной \(t\), измеряемыми в сантиметрах. Плотность алюминия принимается равной
\[
\rho=2.7\ \text{г}/\text{см}^3.
\]
Модель возвращает площадь лицевой поверхности \(A\), объём \(V\) и массу \(m\):
\[
A=wh,
\qquad
V=wht,
\qquad
m=\rho wht.
\]

Для рабочей точки
\[
a=(w,h,t)=(20,10,0.2)
\]
и изменения параметров
\[
\Delta x=(\Delta w,\Delta h,\Delta t)
=(0.1,-0.05,0.01)
\]
построить матрицу Якоби, получить линейный прогноз изменений \(A,V,m\) и сравнить его с точным пересчётом модели.
\end{problem}

\begin{proof}[Решение]
Входной вектор равен
\[
x=(w,h,t)\in\R^3.
\]
Первый столбец будущей матрицы будет соответствовать ширине \(w\), второй --- высоте \(h\), третий --- толщине \(t\).

Выход модели равен
\[
F(w,h,t)
=
\begin{pmatrix}
A(w,h,t)\\
V(w,h,t)\\
m(w,h,t)
\end{pmatrix}
=
\begin{pmatrix}
wh\\
wht\\
\rho wht
\end{pmatrix}
\in\R^3.
\]
Первая строка якобиана должна описывать площадь, вторая --- объём, третья --- массу.

Для площади \(A=wh\) вычисляем чувствительности по каждому входу:
\[
\pdv{A}{w}=h,
\qquad
\pdv{A}{h}=w,
\qquad
\pdv{A}{t}=0.
\]
Нулевая последняя производная появляется потому, что лицевая площадь в принятой модели не зависит от толщины.

Для объёма \(V=wht\):
\[
\pdv{V}{w}=ht,
\qquad
\pdv{V}{h}=wt,
\qquad
\pdv{V}{t}=wh.
\]

Для массы \(m=\rho wht\), где \(\rho\) является постоянной:
\[
\pdv{m}{w}=\rho ht,
\qquad
\pdv{m}{h}=\rho wt,
\qquad
\pdv{m}{t}=\rho wh.
\]

Размещаем найденные производные по правилу «выход --- строка, вход --- столбец»:
\begin{equation}
\Jac{F}(w,h,t)
=
\begin{pmatrix}
h & w & 0\\
ht & wt & wh\\
\rho ht & \rho wt & \rho wh
\end{pmatrix}.
\label{eq:plate-symbolic-jacobian}
\end{equation}

Теперь подставляем рабочую точку \((20,10,0.2)\) и \(\rho=2.7\):
\[
\Jac{F}(a)
=
\begin{pmatrix}
10 & 20 & 0\\
10\cdot0.2 & 20\cdot0.2 & 20\cdot10\\
2.7\cdot10\cdot0.2 & 2.7\cdot20\cdot0.2 & 2.7\cdot20\cdot10
\end{pmatrix}
=
\begin{pmatrix}
10 & 20 & 0\\
2 & 4 & 200\\
5.4 & 10.8 & 540
\end{pmatrix}.
\]

Единицы коэффициентов зависят от строки и столбца. Например,
\[
\pdv{V}{t}(a)=200\ \text{см}^2
\]
означает: увеличение толщины на один сантиметр при фиксированных \(w,h\) увеличивает объём приблизительно на \(200\ \text{см}^3\). Аналогично,
\[
\pdv{m}{t}(a)=540\ \text{г}/\text{см}.
\]

Линейный прогноз изменения выходов равен
\[
\Delta F
\approx
\Jac{F}(a)\Delta x.
\]
Подставляем конкретный вектор изменений:
\[
\Delta F
\approx
\begin{pmatrix}
10 & 20 & 0\\
2 & 4 & 200\\
5.4 & 10.8 & 540
\end{pmatrix}
\begin{pmatrix}
0.1\\
-0.05\\
0.01
\end{pmatrix}.
\]
Первая компонента:
\[
\Delta A_{\mathrm{lin}}
=
10\cdot0.1+20\cdot(-0.05)+0\cdot0.01
=0\ \text{см}^2.
\]
Рост ширины и уменьшение высоты в первом приближении взаимно компенсируются.

Вторая компонента:
\[
\Delta V_{\mathrm{lin}}
=
2\cdot0.1+4\cdot(-0.05)+200\cdot0.01
=2\ \text{см}^3.
\]

Третья компонента:
\[
\Delta m_{\mathrm{lin}}
=
5.4\cdot0.1+10.8\cdot(-0.05)+540\cdot0.01
=5.4\ \text{г}.
\]
Итак,
\[
\Delta F_{\mathrm{lin}}
=
\begin{pmatrix}
0\ \text{см}^2\\
2\ \text{см}^3\\
5.4\ \text{г}
\end{pmatrix}.
\]

Сравним с точным пересчётом. Исходные значения равны
\[
A_0=20\cdot10=200\ \text{см}^2,
\]
\[
V_0=20\cdot10\cdot0.2=40\ \text{см}^3,
\]
\[
m_0=2.7\cdot40=108\ \text{г}.
\]
Новая точка параметров:
\[
a+\Delta x=(20.1,9.95,0.21).
\]
Точные новые значения:
\[
A_1=20.1\cdot9.95=199.995\ \text{см}^2,
\]
\[
V_1=20.1\cdot9.95\cdot0.21
=41.99895\ \text{см}^3,
\]
\[
m_1=2.7\cdot41.99895
=113.397165\ \text{г}.
\]
Следовательно,
\[
\Delta A_{\mathrm{exact}}
=-0.005\ \text{см}^2,
\]
\[
\Delta V_{\mathrm{exact}}
=1.99895\ \text{см}^3,
\]
\[
\Delta m_{\mathrm{exact}}
=5.397165\ \text{г}.
\]

Линейная модель дала соответственно \(0\), \(2\) и \(5.4\). Различия равны
\[
-0.005\ \text{см}^2,
\qquad
-0.00105\ \text{см}^3,
\qquad
-0.002835\ \text{г}.
\]
Они возникают из произведений двух и трёх малых приращений, которые дифференциал не учитывает. Для данного возмущения эти поправки малы, поэтому якобиан корректно описывает локальную чувствительность модели.
\end{proof}


\section{Правило цепочки и распространение чувствительности}
\label{sec:chain-rule-sensitivity}

Во многих моделях результат получается не одной формулой, а последовательностью этапов. Сначала исходные параметры преобразуются во внутренние характеристики, затем эти характеристики используются для вычисления итоговых величин. Правило цепочки показывает, как объединять локальные линейные модели отдельных этапов в локальную линейную модель всей системы.

\subsection{Композиция отображений}

Пусть
\[
f\colon \R^m\to\R^n,
\qquad
g\colon \R^n\to\R^k.
\]
Для входа \(x\in\R^m\) сначала вычисляется промежуточный вектор
\[
y=f(x)\in\R^n,
\]
а затем итоговый вектор
\[
z=g(y)=g(f(x))\in\R^k.
\]
Полное отображение имеет вид
\[
h=g\circ f\colon\R^m\to\R^k.
\]

Если вход изменился на малый вектор \(\Delta x\), то первый этап даёт приближение
\[
\Delta y
\approx
\Jac{f}(x)\Delta x.
\]
Второй этап преобразует полученное изменение:
\[
\Delta z
\approx
\Jac{g}(y)\Delta y.
\]
Подстановка первой формулы во вторую приводит к выражению
\[
\Delta z
\approx
\Jac{g}(f(x))\Jac{f}(x)\Delta x.
\]
Правило цепочки утверждает, что произведение двух матриц действительно является якобианом композиции.

\begin{theorem}[правило цепочки]
Пусть отображение \(f\colon\R^m\to\R^n\) дифференцируемо в точке \(a\), а отображение \(g\colon\R^n\to\R^k\) дифференцируемо в точке
\[
b=f(a).
\]
Тогда композиция \(g\circ f\) дифференцируема в точке \(a\), причём
\begin{equation}
\dd(g\circ f)(a)
=
Dg(b)\circ Df(a).
\label{eq:chain-rule-differential}
\end{equation}
\end{theorem}

\begin{proof}
Обозначим
\[
A=Df(a),
\qquad
B=Dg(b).
\]
Дифференцируемость \(f\) означает, что
\[
f(a+h)=b+Ah+r_f(h),
\qquad
\frac{\norm{r_f(h)}}{\norm{h}}\to0.
\]
Для \(g\) имеем
\[
g(b+u)=g(b)+Bu+r_g(u),
\qquad
\frac{\norm{r_g(u)}}{\norm{u}}\to0.
\]
Положим
\[
u=Ah+r_f(h).
\]
Тогда \(u=O(\norm{h})\), поскольку линейное отображение \(A\) ограничено, а \(r_f(h)=o(\norm{h})\). Подставляя выражение для \(f(a+h)\) во второй этап, получаем
\[
(g\circ f)(a+h)
=
g(b)+B\bigl(Ah+r_f(h)\bigr)+r_g(u).
\]
Следовательно,
\[
(g\circ f)(a+h)-(g\circ f)(a)
=
BAh+\bigl(Br_f(h)+r_g(u)\bigr).
\]
Первое слагаемое линейно по \(h\). Кроме того,
\[
Br_f(h)=o(\norm{h}),
\qquad
r_g(u)=o(\norm{u})=o(\norm{h}).
\]
Поэтому остаток в скобках имеет порядок \(o(\norm{h})\), а дифференциал композиции равен \(BA\), то есть композиции \(B\circ A\).
\end{proof}

\subsection{Матричная форма и порядок множителей}

В стандартных базисах дифференциалы задаются матрицами Якоби:
\[
\Jac{f}(a)\in\R^{n\times m},
\qquad
\Jac{g}(b)\in\R^{k\times n}.
\]
Поэтому формула \eqref{eq:chain-rule-differential} принимает вид
\begin{equation}
\Jac{(g\circ f)}(a)
=
\Jac{g}(f(a))\Jac{f}(a).
\label{eq:chain-rule-jacobian}
\end{equation}
Размерности согласуются:
\[
(k\times n)(n\times m)=k\times m.
\]
Полученная матрица переводит изменение исходного входа из \(\R^m\) непосредственно в изменение конечного выхода из \(\R^k\).

Порядок множителей определяется порядком выполнения отображений. Сначала на вектор \(\Delta x\) действует \(\Jac{f}(a)\), а затем на результат действует \(\Jac{g}(b)\):
\[
\Delta x
\xmapsto{\ \Jac{f}(a)\ }
\Delta y
\xmapsto{\ \Jac{g}(b)\ }
\Delta z.
\]
Поэтому в произведении матрица второго этапа стоит слева.

\begin{caution}
В общем случае
\[
\Jac{g}(b)\Jac{f}(a)
\neq
\Jac{f}(a)\Jac{g}(b).
\]
Обратный порядок часто даже не имеет смысла по размерностям. Проверка размеров является быстрым способом обнаружить ошибку.
\end{caution}

Пусть
\[
f=(f_1,\dots,f_n),
\qquad
g=(g_1,\dots,g_k).
\]
Элемент строки \(\ell\) и столбца \(i\) итогового якобиана равен
\begin{equation}
\pdv{(g_\ell\circ f)}{x_i}(a)
=
\sum_{j=1}^{n}
\pdv{g_\ell}{y_j}(b)
\pdv{f_j}{x_i}(a).
\label{eq:chain-rule-coordinate}
\end{equation}
Каждое слагаемое описывает один путь влияния
\[
x_i\longrightarrow y_j\longrightarrow z_\ell.
\]
Производная \(\pdv{f_j}{x_i}\) измеряет влияние входа \(x_i\) на промежуточную величину \(y_j\), а производная \(\pdv{g_\ell}{y_j}\) --- влияние этой промежуточной величины на конечный выход \(z_\ell\). Сумма объединяет все пути через компоненты \(y_1,\dots,y_n\).

\subsection{Скалярный выход и перенос градиента}

Пусть теперь
\[
f\colon\R^m\to\R^n,
\qquad
g\colon\R^n\to\R,
\qquad
q=g\circ f.
\]
Якобиан скалярной функции \(g\) является строкой
\[
\Jac{g}(y)=\grad g(y)^{\trans}.
\]
Из правила цепочки следует
\[
\grad q(a)^{\trans}
=
\grad g(f(a))^{\trans}\Jac{f}(a).
\]
После транспонирования получаем
\begin{equation}
\grad(g\circ f)(a)
=
\Jac{f}(a)^{\trans}\grad g(f(a)).
\label{eq:gradient-chain-rule}
\end{equation}

Формула показывает два направления вычисления. При прямом распространении малое изменение входа умножается на якобианы слева направо по ходу модели. При вычислении градиента скалярного результата вектор чувствительности переносится назад транспонированными якобианами.

Для нескольких последовательных этапов правило применяется повторно. Промежуточные якобианы не обязаны быть квадратными: достаточно совпадения размерности выхода предыдущего этапа с размерностью входа следующего.

\subsection{Изменение функции вдоль траектории}

Пусть состояние системы зависит от времени:
\[
x=x(t)\in\R^m,
\]
а скалярная характеристика состояния задаётся функцией \(F\colon\R^m\to\R\). Наблюдаемая во времени величина имеет вид композиции
\[
\varphi(t)=F(x(t)).
\]
Правило цепочки даёт
\begin{equation}
\dv{}{t}F(x(t))
=
\grad F(x(t))^{\trans}x'(t)
=
\sum_{i=1}^{m}
\pdv{F}{x_i}(x(t))x_i'(t).
\label{eq:trajectory-chain-rule}
\end{equation}

Вектор \(x'(t)\) описывает скорость изменения состояния, а градиент --- чувствительность характеристики к координатам состояния. Их скалярное произведение определяет мгновенную скорость изменения наблюдаемой величины.

Если \(F\colon\R^m\to\R^k\) имеет векторный выход, то
\begin{equation}
\dv{}{t}F(x(t))
=
\Jac{F}(x(t))x'(t).
\label{eq:vector-trajectory-chain-rule}
\end{equation}
Строка \(j\) матрицы Якоби вычисляет скорость изменения \(j\)-й выходной компоненты.

\subsection{Абсолютная и относительная чувствительность}

Для скалярного выхода \(q(x)\) производная \(\pdv{q}{x_i}(a)\) является локальной абсолютной чувствительностью: она имеет единицы «выход на единицу входа» и показывает изменение \(q\) при малом изменении одного параметра. Если \(a_i\neq0\) и \(q(a)\neq0\), безразмерная относительная чувствительность определяется формулой
\begin{equation}
E_i(a)=\frac{a_i}{q(a)}\pdv{q}{x_i}(a),
\qquad
\frac{\Delta q}{q(a)}
\approx
\sum_{i=1}^{m}E_i(a)\frac{\Delta x_i}{a_i}.
\label{eq:relative-sensitivity}
\end{equation}
Например, \(E_i=2\) означает, что малое увеличение \(x_i\) на один процент даёт приблизительно двухпроцентное увеличение выхода при фиксированных остальных параметрах. При нулевом входе или выходе используют абсолютную форму через дифференциал.

\begin{problem}[Задача по материалу В.~А.~Зорича: поле вдоль траектории]
Нормированное давление в установившемся потоке задано функцией
\[
p(x,y,z)=x^2y+e^z.
\]
Датчик движется по траектории
\[
r(t)=(1+t,\,2-t,\,t^2).
\]
Найти скорость изменения показания датчика в момент \(t=0\) двумя способами: по правилу цепочки и прямым дифференцированием композиции.
\end{problem}

\begin{proof}[Решение]
Промежуточное отображение \(r\colon\R\to\R^3\) имеет производную
\[
r'(t)
=
\begin{pmatrix}
1\\
-1\\
2t
\end{pmatrix}.
\]
В момент \(t=0\)
\[
r(0)=(1,2,0),
\qquad
r'(0)=
\begin{pmatrix}
1\\
-1\\
0
\end{pmatrix}.
\]

Градиент давления равен
\[
\grad p(x,y,z)
=
\begin{pmatrix}
2xy\\
x^2\\
e^z
\end{pmatrix}.
\]
В точке \(r(0)
=(1,2,0)\) получаем
\[
\grad p(r(0))
=
\begin{pmatrix}
4\\
1\\
1
\end{pmatrix}.
\]
По формуле \eqref{eq:trajectory-chain-rule}
\[
\dv{}{t}p(r(t))\bigg|_{t=0}
=
\grad p(r(0))^{\trans}r'(0)
=
\begin{pmatrix}4&1&1\end{pmatrix}
\begin{pmatrix}1\\-1\\0\end{pmatrix}
=3.
\]

Проверим результат прямой подстановкой:
\[
p(r(t))
=(1+t)^2(2-t)+e^{t^2}.
\]
Дифференцируем:
\[
\dv{}{t}p(r(t))
=
2(1+t)(2-t)-(1+t)^2+2t e^{t^2}.
\]
При \(t=0\)
\[
\dv{}{t}p(r(t))\bigg|_{t=0}
=4-1+0=3.
\]
Оба способа дали один результат. Первый способ разделяет геометрию движения и свойства поля, поэтому удобен, когда траектория или поле меняются независимо друг от друга.
\end{proof}

\begin{problem}[Практическая задача: двухэтапная модель цилиндрической детали]
Закрытая цилиндрическая деталь задаётся радиусом \(r\) и длиной \(L\), измеряемыми в сантиметрах. Первый этап модели вычисляет полную площадь поверхности и объём:
\[
f(r,L)
=
\begin{pmatrix}
A(r,L)\\
V(r,L)
\end{pmatrix}
=
\begin{pmatrix}
2\pi rL+2\pi r^2\\
\pi r^2L
\end{pmatrix}.
\]
Второй этап вычисляет стоимость материала, массу и среднюю нагрузку на единицу площади:
\[
g(A,V)
=
\begin{pmatrix}
C\\
m\\
q
\end{pmatrix}
=
\begin{pmatrix}
c_AA+c_VV\\
\rho V\\
P/A
\end{pmatrix}.
\]
Пусть
\[
r=2\ \text{см},
\qquad
L=10\ \text{см},
\]
\[
c_A=0.05\ \text{руб}/\text{см}^2,
\qquad
c_V=0.02\ \text{руб}/\text{см}^3,
\]
\[
\rho=2.7\ \text{г}/\text{см}^3,
\qquad
P=300\ \text{Н}.
\]
Требуется построить полный якобиан композиции \(g\circ f\), оценить изменение выходов при
\[
\Delta r=0.02\ \text{см},
\qquad
\Delta L=-0.10\ \text{см},
\]
сравнить прогноз с точным пересчётом и найти относительные чувствительности массы.
\end{problem}

\begin{proof}[Решение]
Исходный вектор и промежуточный вектор имеют размерности
\[
x=(r,L)\in\R^2,
\qquad
y=(A,V)\in\R^2,
\]
а итоговый выход
\[
z=(C,m,q)\in\R^3.
\]
Следовательно,
\[
\Jac{f}(x)\in\R^{2\times2},
\qquad
\Jac{g}(y)\in\R^{3\times2},
\]
и полный якобиан должен иметь размер \(3\times2\).

Для геометрического этапа
\[
\pdv{A}{r}=2\pi L+4\pi r,
\qquad
\pdv{A}{L}=2\pi r,
\]
\[
\pdv{V}{r}=2\pi rL,
\qquad
\pdv{V}{L}=\pi r^2.
\]
Поэтому
\begin{equation}
\Jac{f}(r,L)
=
\begin{pmatrix}
2\pi L+4\pi r & 2\pi r\\
2\pi rL & \pi r^2
\end{pmatrix}.
\label{eq:cylinder-geometry-jacobian}
\end{equation}
В рабочей точке \((2,10)\)
\[
\Jac{f}(2,10)
=
\begin{pmatrix}
28\pi & 4\pi\\
40\pi & 4\pi
\end{pmatrix}.
\]
Первая строка описывает изменение площади, вторая --- изменение объёма. Первый столбец относится к радиусу, второй --- к длине.

Промежуточные значения равны
\[
A_0=2\pi\cdot2\cdot10+2\pi\cdot2^2=48\pi\ \text{см}^2,
\]
\[
V_0=\pi\cdot2^2\cdot10=40\pi\ \text{см}^3.
\]
Численно
\[
A_0\approx150.796\ \text{см}^2,
\qquad
V_0\approx125.664\ \text{см}^3.
\]

Для второго этапа
\[
\Jac{g}(A,V)
=
\begin{pmatrix}
 c_A & c_V\\
 0 & \rho\\
 -P/A^2 & 0
\end{pmatrix}.
\]
Строки относятся к стоимости, массе и нагрузке, а столбцы --- к площади и объёму. В точке \((A_0,V_0)\)
\[
\Jac{g}(A_0,V_0)
=
\begin{pmatrix}
0.05 & 0.02\\
0 & 2.7\\
-\dfrac{300}{(48\pi)^2} & 0
\end{pmatrix}.
\]

По правилу цепочки
\[
\Jac{(g\circ f)}(2,10)
=
\Jac{g}(A_0,V_0)\Jac{f}(2,10).
\]
Выполняем умножение строк на столбцы:
\begin{equation}
\Jac{(g\circ f)}(2,10)
=
\begin{pmatrix}
2.2\pi & 0.28\pi\\
108\pi & 10.8\pi\\
-\dfrac{175}{48\pi} & -\dfrac{25}{48\pi}
\end{pmatrix}.
\label{eq:cylinder-total-jacobian}
\end{equation}
Например, первый элемент первой строки получен как
\[
0.05\cdot28\pi+0.02\cdot40\pi=2.2\pi.
\]
Это сумма двух путей влияния радиуса на стоимость: через площадь и через объём.

Численная форма полного якобиана:
\[
\Jac{(g\circ f)}(2,10)
\approx
\begin{pmatrix}
6.9115 & 0.8796\\
339.2920 & 33.9292\\
-1.1605 & -0.1658
\end{pmatrix}.
\]
Единицы по строкам различаются: первая строка измеряется в рублях на сантиметр, вторая --- в граммах на сантиметр, третья --- в \(\text{Н}/\text{см}^3\).

Линейный прогноз равен
\[
\Delta z
\approx
\Jac{(g\circ f)}(2,10)
\begin{pmatrix}
0.02\\
-0.10
\end{pmatrix}.
\]
Покомпонентно получаем
\begin{align*}
\Delta C_{\mathrm{lin}}
&=2.2\pi\cdot0.02+0.28\pi\cdot(-0.10)
=0.016\pi\approx0.05027\ \text{руб},\\
\Delta m_{\mathrm{lin}}
&=108\pi\cdot0.02+10.8\pi\cdot(-0.10)
=1.08\pi\approx3.39292\ \text{г},\\
\Delta q_{\mathrm{lin}}
&=-\frac{175}{48\pi}\cdot0.02-\frac{25}{48\pi}\cdot(-0.10)
=-\frac{1}{48\pi}\approx-0.006631\ \text{Н}/\text{см}^2.
\end{align*}

Проверим результат точным пересчётом. Новые размеры:
\[
r_1=2.02\ \text{см},
\qquad
L_1=9.90\ \text{см}.
\]
Тогда
\[
A_1
=2\pi r_1L_1+2\pi r_1^2
=48.1568\pi
\approx151.28905\ \text{см}^2,
\]
\[
V_1
=\pi r_1^2L_1
=40.39596\pi
\approx126.90765\ \text{см}^3.
\]
Исходные значения и точные изменения равны
\begin{align*}
C_0&=3.2\pi\approx10.05310\ \text{руб},
&\Delta C_{\mathrm{exact}}&\approx0.049509\ \text{руб},\\
m_0&=108\pi\approx339.29201\ \text{г},
&\Delta m_{\mathrm{exact}}&\approx3.358652\ \text{г},\\
q_0&=\frac{300}{48\pi}\approx1.989437\ \text{Н}/\text{см}^2,
&\Delta q_{\mathrm{exact}}&\approx-0.006478\ \text{Н}/\text{см}^2.
\end{align*}
Разности между точными и линейными изменениями равны приблизительно
\[
-0.000756\ \text{руб},
\qquad
-0.034268\ \text{г},
\qquad
0.000154\ \text{Н}/\text{см}^2.
\]
Они малы по сравнению с самими выходами, поэтому локальная линейная модель подходит для заданного возмущения.

Наконец, масса имеет явную формулу
\[
m(r,L)=\rho\pi r^2L.
\]
Её относительные чувствительности равны
\[
E_r
=
\frac{r}{m}\pdv{m}{r}
=
\frac{r}{\rho\pi r^2L}\,2\rho\pi rL
=2,
\]
\[
E_L
=
\frac{L}{m}\pdv{m}{L}
=
\frac{L}{\rho\pi r^2L}\,\rho\pi r^2
=1.
\]
Следовательно, увеличение радиуса на один процент при фиксированной длине увеличивает массу приблизительно на два процента, а увеличение длины на один процент --- приблизительно на один процент. Этот вывод не зависит от конкретной рабочей точки, пока сохраняется цилиндрическая модель.
\end{proof}


\section{Второй порядок и ограничения}
\label{sec:second-order-implicit-constrained}

Дифференциал описывает главный линейный эффект малого изменения входа. Вблизи стационарной точки этот эффект исчезает, поскольку градиент равен нулю. Тогда поведение функции определяет следующий, квадратичный член. Тот же аппарат позволяет понять локальную обратимость системы координат, вычислять чувствительность переменной, заданной неявным уравнением, и искать экстремум при наличии связей.

\subsection{Матрица Гессе и квадратичная модель}

Пусть функция
\[
f\colon G\subseteq\R^n\to\R
\]
имеет в окрестности точки \(a\in G\) непрерывные частные производные второго порядка.

\begin{definition}
Матрицей Гессе функции \(f\) в точке \(a\) называется квадратная матрица
\[
\Hessian{f}(a)
=
\begin{pmatrix}
\dfrac{\partial^2 f}{\partial x_1^2}(a)
& \cdots &
\dfrac{\partial^2 f}{\partial x_1\partial x_n}(a)\\[1.1ex]
\vdots & \ddots & \vdots\\[0.6ex]
\dfrac{\partial^2 f}{\partial x_n\partial x_1}(a)
& \cdots &
\dfrac{\partial^2 f}{\partial x_n^2}(a)
\end{pmatrix}.
\]
\end{definition}

Элемент с индексами \((i,j)\) измеряет, как изменяется производная по \(x_i\) при изменении \(x_j\). Диагональные элементы описывают кривизну вдоль отдельных координат, а внедиагональные --- взаимодействие двух переменных. Если вторые производные непрерывны, то смешанные производные совпадают, поэтому
\[
\Hessian{f}(a)^{\trans}=\Hessian{f}(a).
\]

Для приращения \(h=(h_1,\dots,h_n)\) второй дифференциал на паре одинаковых векторов записывается как квадратичная форма
\[
\dd^2 f(a)[h,h]
=
\sum_{i,j=1}^n
\frac{\partial^2 f}{\partial x_i\partial x_j}(a)h_i h_j
=
h^{\trans}\Hessian{f}(a)h.
\]
Размерность каждого слагаемого совпадает с размерностью выхода \(f\). Например, если \(f\) измеряется в рублях, а \(x_i\) и \(x_j\) --- в сантиметрах, то соответствующий элемент Гессе имеет единицы рублей на квадратный сантиметр.

\begin{theorem}[формула Тейлора второго порядка]
Если \(f\in C^2\) в окрестности точки \(a\), то при \(h\to0\)
\begin{equation}
 f(a+h)
 =
 f(a)
 +\grad f(a)^{\trans}h
 +\frac12 h^{\trans}\Hessian{f}(a)h
 +o\bigl(\norm{h}^2\bigr).
\label{eq:taylor-second-order}
\end{equation}
\end{theorem}

Формула получается применением одномерной формулы Тейлора к функции
\[
\varphi(t)=f(a+th),
\qquad 0\leq t\leq1.
\]
В точке \(t=0\)
\[
\varphi'(0)=\grad f(a)^{\trans}h,
\qquad
\varphi''(0)=h^{\trans}\Hessian{f}(a)h.
\]
Следовательно, многомерная формула представляет собой обычное разложение вдоль прямого отрезка от \(a\) к \(a+h\).

\subsection{Стационарные точки и достаточный признак экстремума}

\begin{definition}
Точка \(a\) называется стационарной для \(f\), если
\[
\grad f(a)=0.
\]
\end{definition}

Во внутренней точке локального экстремума дифференцируемой функции градиент обязан обращаться в нуль. Обратное утверждение неверно: стационарная точка может быть минимумом, максимумом или седловой точкой.

В стационарной точке линейный член в \eqref{eq:taylor-second-order} исчезает:
\[
f(a+h)-f(a)
=
\frac12 h^{\trans}\Hessian{f}(a)h
+o\bigl(\norm{h}^2\bigr).
\]
Поэтому знак квадратичной формы определяет локальное поведение функции.

\begin{criterion}[критерий второго порядка]
Пусть \(a\) --- стационарная точка функции \(f\in C^2\). Тогда:
Если \(h^{\trans}\Hessian{f}(a)h>0\) при каждом \(h\neq0\), точка \(a\) является строгим локальным минимумом; если эта форма всюду отрицательна --- строгим локальным максимумом. При значениях обоих знаков точка экстремумом не является.
Если форма лишь полуопределена, критерий не даёт ответа: необходимо исследовать члены более высокого порядка или саму функцию.
\end{criterion}

Для функции двух переменных положим
\[
D(a)
=
\det\Hessian{f}(a)
=
f_{xx}(a)f_{yy}(a)-f_{xy}(a)^2.
\]
Тогда при \(D(a)>0\) знак \(f_{xx}(a)\) различает минимум и максимум, а при \(D(a)<0\) точка седловая. Случай \(D(a)=0\) остаётся неопределённым.

\subsection{Неявная функция как зависимость, заданная уравнением}

Во многих моделях зависимая величина не выражена готовой формулой. Вместо \(y=f(x)\) задаётся уравнение
\[
F(x,y)=0.
\]
Теорема о неявной функции указывает, когда это уравнение всё же определяет локальную функцию.

\begin{theorem}[неявная функция, рабочая форма]
Пусть \(F\colon\R^m\times\R^k\to\R^k\) непрерывно дифференцируема в окрестности точки \((x_0,y_0)\), причём
\[
F(x_0,y_0)=0
\]
и квадратная матрица частных производных по зависимым переменным
\[
F_y(x_0,y_0)
\]
обратима. Тогда в некоторой окрестности \(x_0\) уравнение \(F(x,y)=0\) однозначно задаёт функцию \(y=\varphi(x)\), причём
\begin{equation}
\Jac{\varphi}(x)
=
-\bigl[F_y(x,\varphi(x))\bigr]^{-1}F_x(x,\varphi(x)).
\label{eq:implicit-derivative}
\end{equation}
\end{theorem}

В скалярном случае \(F\colon\R^2\to\R\) условие принимает вид
\[
F_y(x_0,y_0)\neq0,
\]
а формула производной упрощается:
\begin{equation}
\varphi'(x)
=-\frac{F_x(x,\varphi(x))}{F_y(x,\varphi(x))}.
\label{eq:implicit-scalar-derivative}
\end{equation}
Знаменатель показывает, почему требуется условие \(F_y\neq0\): если изменение \(y\) в первом порядке не меняет левую часть уравнения, то локально восстановить \(y\) по \(x\) может быть невозможно.

\subsection{Локальная обратимость отображения}

\begin{theorem}[обратная функция, рабочая форма]
Пусть \(f\colon G\subseteq\R^n\to\R^n\) непрерывно дифференцируема и
\[
\det\Jac{f}(a)\neq0.
\]
Тогда в некоторой окрестности точки \(a\) отображение имеет дифференцируемое обратное. Если \(b=f(a)\), то
\begin{equation}
\Jac{(f^{-1})}(b)
=
\bigl[\Jac{f}(a)\bigr]^{-1}.
\label{eq:inverse-function-derivative}
\end{equation}
\end{theorem}

Условие не утверждает глобальной взаимной однозначности. Оно гарантирует только локальную разрешимость системы \(y=f(x)\) относительно \(x\). Например, переход
\[
(r,\theta)
\longmapsto
(x,y)
=
(r\cos\theta,r\sin\theta)
\]
имеет определитель Якоби \(r\). При \(r>0\) координаты локально обратимы, а при \(r=0\) угол восстановить нельзя.

\subsection{Условный экстремум и множители Лагранжа}

Пусть требуется найти экстремум функции \(f(x)\) не во всём пространстве, а на множестве, заданном связями
\[
F_1(x)=0,
\quad\dots,\quad
F_m(x)=0.
\]
Допустимое приращение \(h\) должно в первом порядке сохранять каждую связь:
\[
\grad F_i(a)^{\trans}h=0,
\qquad i=1,\dots,m.
\]
В точке условного экстремума также должно выполняться
\[
\grad f(a)^{\trans}h=0
\]
для всех допустимых направлений. Поэтому градиент цели лежит в линейной оболочке градиентов ограничений.

\begin{theorem}[необходимое условие Лагранжа]
Пусть \(a\) является локальным экстремумом \(f\) при связях \(F_i(x)=0\), а векторы
\[
\grad F_1(a),\dots,\grad F_m(a)
\]
линейно независимы. Тогда существуют числа \(\lambda_1,\dots,\lambda_m\), для которых
\begin{equation}
\grad f(a)
=
\sum_{i=1}^m\lambda_i\grad F_i(a).
\label{eq:lagrange-gradient}
\end{equation}
\end{theorem}

Вводя функцию Лагранжа
\[
L(x,\lambda)
=
f(x)-\sum_{i=1}^m\lambda_iF_i(x),
\]
получаем систему
\[
\grad_x L(x,\lambda)=0,
\qquad
F_1(x)=\cdots=F_m(x)=0.
\]
Это лишь система кандидатов. После её решения необходимо сравнить значения цели, проверить границы области и, когда требуется локальная классификация, исследовать второй порядок вдоль допустимых направлений.

\begin{problem}[Задача по материалу В.~А.~Зорича: квадратичная энергия на сфере]
Пусть
\[
E(x,y)=3x^2+2xy+3y^2
\]
--- энергия двух связанных компонент, а условие
\[
x^2+y^2=1
\]
фиксирует общий уровень состояния. Найти минимальное и максимальное значения энергии и соответствующие состояния.
\end{problem}

\begin{proof}[Решение]
Запишем связь как
\[
F(x,y)=x^2+y^2-1=0
\]
и введём функцию Лагранжа
\[
L(x,y,\lambda)
=
3x^2+2xy+3y^2-\lambda(x^2+y^2-1).
\]
Уравнения стационарности имеют вид
\[
\pdv{L}{x}=6x+2y-2\lambda x=0,
\]
\[
\pdv{L}{y}=2x+6y-2\lambda y=0,
\]
\[
\pdv{L}{\lambda}=-(x^2+y^2-1)=0.
\]
Первые два уравнения после деления на \(2\) записываются матрично:
\[
\begin{pmatrix}
3&1\\
1&3
\end{pmatrix}
\begin{pmatrix}
x\\y
\end{pmatrix}
=
\lambda
\begin{pmatrix}
x\\y
\end{pmatrix}.
\]
Таким образом, допустимое состояние должно быть собственным вектором симметрической матрицы энергии. Её характеристическое уравнение
\[
\det
\begin{pmatrix}
3-\lambda&1\\
1&3-\lambda
\end{pmatrix}
=(3-\lambda)^2-1=0
\]
даёт
\[
\lambda_1=2,
\qquad
\lambda_2=4.
\]
При \(\lambda_1=2\) имеем \(x+y=0\). С учётом единичной нормы получаем два состояния
\[
\frac{1}{\sqrt2}(1,-1),
\qquad
\frac{1}{\sqrt2}(-1,1).
\]
В них
\[
E=2.
\]
При \(\lambda_2=4\) имеем \(x-y=0\), поэтому
\[
\frac{1}{\sqrt2}(1,1),
\qquad
\frac{1}{\sqrt2}(-1,-1),
\]
и
\[
E=4.
\]
Сфера \(x^2+y^2=1\) компактна, а энергия непрерывна, поэтому найденные наименьшее и наибольшее значения действительно являются глобальными. Противофазный режим \(x=-y\) минимизирует энергию, а синфазный режим \(x=y\) максимизирует её.
\end{proof}

\begin{problem}[Практическая задача: цилиндр минимальной материалоёмкости]
Требуется изготовить закрытый цилиндрический резервуар объёма
\[
V_0=1000\ \text{см}^3.
\]
Радиус \(r>0\) и высота \(h>0\) можно выбирать. Толщина материала постоянна, поэтому материалоёмкость пропорциональна площади полной поверхности
\[
S(r,h)=2\pi r^2+2\pi rh.
\]
Нужно:

Найти размеры, минимизирующие \(S\) при условии \(\pi r^2h=V_0\), проверить минимум, вычислить размеры и площадь, определить их чувствительность к \(V_0\) и оценить рост площади при малом отклонении радиуса с сохранением объёма.
\end{problem}

\begin{proof}[Решение]
Целевая функция и связь равны
\[
S(r,h)=2\pi r^2+2\pi rh,
\qquad
F(r,h)=\pi r^2h-V_0=0.
\]
Их градиенты:
\[
\grad S(r,h)
=
\begin{pmatrix}
4\pi r+2\pi h\\
2\pi r
\end{pmatrix},
\qquad
\grad F(r,h)
=
\begin{pmatrix}
2\pi rh\\
\pi r^2
\end{pmatrix}.
\]
Условие Лагранжа \(\grad S=\lambda\grad F\) даёт систему
\[
4\pi r+2\pi h
=
\lambda\,2\pi rh,
\]
\[
2\pi r
=
\lambda\,\pi r^2,
\]
\[
\pi r^2h=V_0.
\]
Поскольку \(r>0\), из второго уравнения
\[
\lambda=\frac{2}{r}.
\]
Подставляем это значение в первое уравнение:
\[
4\pi r+2\pi h
=
\frac{2}{r}\,2\pi rh
=4\pi h.
\]
После сокращения на \(2\pi\)
\[
2r+h=2h,
\qquad
h=2r.
\]
Итак, оптимальная высота равна диаметру. Подстановка в ограничение даёт
\[
2\pi r^3=V_0,
\]
откуда
\begin{equation}
r_*(V_0)
=
\left(\frac{V_0}{2\pi}\right)^{1/3},
\qquad
h_*(V_0)
=
2\left(\frac{V_0}{2\pi}\right)^{1/3}.
\label{eq:optimal-cylinder-dimensions}
\end{equation}

Проверим тип точки без громоздкого условного Гессе. Из связи выражаем высоту:
\[
h(r)=\frac{V_0}{\pi r^2}.
\]
Это допустимо при \(r>0\), поскольку
\[
F_h(r,h)=\pi r^2\neq0.
\]
Тем самым используется теорема о неявной функции. После подстановки площадь становится функцией одной переменной:
\[
\widetilde S(r)
=
2\pi r^2
+2\pi r\frac{V_0}{\pi r^2}
=
2\pi r^2+\frac{2V_0}{r}.
\]
Её производные равны
\[
\widetilde S'(r)
=
4\pi r-\frac{2V_0}{r^2},
\]
\[
\widetilde S''(r)
=
4\pi+\frac{4V_0}{r^3}.
\]
В точке \(r=r_*\) выполнено \(V_0=2\pi r_*^3\), поэтому
\[
\widetilde S''(r_*)
=
4\pi+8\pi
=12\pi>0.
\]
Следовательно, стационарная точка является строгим локальным минимумом. Кроме того,
\[
\widetilde S(r)\to+\infty
\quad\text{при}\quad
r\to0+,
\qquad
\widetilde S(r)\to+\infty
\quad\text{при}\quad
r\to+\infty,
\]
поэтому минимум глобальный.

Для \(V_0=1000\ \text{см}^3\)
\[
r_*
=
\left(\frac{1000}{2\pi}\right)^{1/3}
\approx5.4193\ \text{см},
\]
\[
h_*=2r_*
\approx10.8385\ \text{см}.
\]
Минимальная площадь
\[
S_*
=2\pi r_*^2+2\pi r_*h_*
=6\pi r_*^2
\approx553.5810\ \text{см}^2.
\]

Рассмотрим изменение требуемого объёма. Оптимальный радиус неявно задаётся уравнением
\[
G(r,V)=2\pi r^3-V=0.
\]
Поскольку
\[
G_r(r,V)=6\pi r^2>0,
\]
радиус является гладкой функцией объёма. По формуле \eqref{eq:implicit-scalar-derivative}
\[
\dv{r_*}{V}
=-\frac{G_V}{G_r}
=
\frac{1}{6\pi r_*^2}.
\]
При \(V_0=1000\ \text{см}^3\)
\[
\dv{r_*}{V}
\approx0.0018064\ \frac{\text{см}}{\text{см}^3}.
\]
Из явной формулы \(r_*\propto V^{1/3}\) следует относительная чувствительность
\[
\frac{V}{r_*}\dv{r_*}{V}
=\frac13.
\]
Поэтому увеличение требуемого объёма на \(1\%\) увеличивает оптимальные радиус и высоту приблизительно на \(1/3\%\).

Наконец, пусть при сохранении объёма радиус изменён на \(\Delta r\). В оптимальной точке \(\widetilde S'(r_*)=0\), поэтому линейный вклад исчезает. Формула Тейлора второго порядка даёт
\[
\Delta S
=
\widetilde S(r_*+\Delta r)-\widetilde S(r_*)
\approx
\frac12\widetilde S''(r_*)(\Delta r)^2
=6\pi(\Delta r)^2.
\]
Например, при \(\Delta r=0.10\ \text{см}\)
\[
\Delta S
\approx
6\pi\cdot0.10^2
\approx0.1885\ \text{см}^2.
\]
Первый порядок здесь не видит потери материалоёмкости, а второй порядок даёт её главный ненулевой вклад.
\end{proof}


\section{Кратные интегралы}
\label{sec:multiple-integrals}

Дифференциал описывает локальное изменение функции вблизи одной точки. Кратный интеграл решает другую задачу: он суммирует локальные вклады по всей области. Если область разбить на малые ячейки, в каждой ячейке значение функции умножается на площадь или объём ячейки, а затем все произведения складываются. Предел таких сумм и является интегралом.

В этом разделе рассматривается интеграл Римана в конечномерном пространстве. Основное внимание уделено вычислительной стороне: как задать область, откуда берутся пределы повторного интеграла и как интерпретировать результат в задачах о массе, среднем значении и центре масс.

\subsection{Координатный параллелепипед и его объём}

\begin{definition}
Координатным параллелепипедом в $\R^n$ называется множество
\[
I=[a_1,b_1]\times\cdots\times[a_n,b_n]
=
\{x\in\R^n\mid a_i\leq x_i\leq b_i,
\ i=1,\dots,n\}.
\]
Его $n$-мерный объём определяется формулой
\begin{equation}
\abs{I}
=
\prod_{i=1}^{n}(b_i-a_i).
\label{eq:n-box-volume}
\end{equation}
\end{definition}

При $n=1$ это длина отрезка, при $n=2$ --- площадь прямоугольника, при $n=3$ --- объём прямоугольного параллелепипеда. Формула \eqref{eq:n-box-volume} показывает, что единица измерения объёма получается перемножением единиц всех координат. Например, если $x$ и $y$ измеряются в метрах, то $\dd x\,\dd y$ имеет размерность $\text{м}^2$.

Пусть параллелепипед $I$ разбит на меньшие координатные параллелепипеды
\[
I_1,\dots,I_N,
\]
не имеющие общих внутренних точек. В каждой ячейке $I_k$ выбирается отмеченная точка $\xi_k\in I_k$. Диаметр ячейки равен наибольшему расстоянию между двумя её точками, а параметр разбиения определяется как
\[
\lambda(P)=\max_{1\leq k\leq N}\operatorname{diam}I_k.
\]
Условие $\lambda(P)\to0$ означает, что все ячейки одновременно становятся малыми.

\subsection{Интегральная сумма и интеграл Римана}

\begin{definition}
Для функции $f\colon I\to\R$ сумма
\begin{equation}
\sigma(f;P,\xi)
=
\sum_{k=1}^{N}f(\xi_k)\abs{I_k}
\label{eq:multiple-riemann-sum}
\end{equation}
называется интегральной суммой Римана.
\end{definition}

Каждое слагаемое имеет прямой смысл. Значение $f(\xi_k)$ характеризует интенсивность величины в выбранной точке ячейки, а $\abs{I_k}$ --- размер самой ячейки. Произведение приближённо равно полному вкладу ячейки.

\begin{definition}
Если существует предел интегральных сумм, не зависящий от способа разбиения и выбора отмеченных точек, то функция $f$ называется интегрируемой по Риману на $I$, а число
\begin{equation}
\int_I f(x)\,\dd x
=
\lim_{\lambda(P)\to0}
\sum_{k=1}^{N}f(\xi_k)\abs{I_k}
\label{eq:multiple-riemann-integral}
\end{equation}
называется кратным интегралом функции $f$ по $I$.
\end{definition}

Для $n=2$ используются записи
\[
\iint_I f(x,y)\,\dd x\,\dd y,
\]
а для $n=3$ ---
\[
\iiint_I f(x,y,z)\,\dd x\,\dd y\,\dd z.
\]
Порядок дифференциалов в такой записи пока не означает порядок вычислений: это единый интеграл по всей области. Порядок появится после применения теоремы Фубини.

\begin{theorem}[достаточное условие интегрируемости]
Всякая функция, непрерывная на координатном параллелепипеде $I\subset\R^n$, интегрируема по Риману на $I$.
\end{theorem}

Непрерывность на компакте обеспечивает равномерную непрерывность. Поэтому при достаточно малых ячейках значения функции внутри каждой ячейки мало отличаются друг от друга, и выбор отмеченной точки перестаёт существенно влиять на сумму.

\begin{remark}
Интегрируемыми бывают и разрывные функции. Для курса вычислительного многомерного анализа достаточно помнить практический критерий: ограниченная функция с разрывами только на конечном числе гладких кривых или поверхностей интегрируема по Риману на ограниченной области.
\end{remark}

\subsection{Интеграл по области}

Пусть $D\subset\R^n$ --- ограниченное множество, содержащееся в некотором координатном параллелепипеде $I$. Введём характеристическую функцию
\[
\mathbf 1_D(x)
=
\begin{cases}
1,&x\in D,\\
0,&x\notin D.
\end{cases}
\]

\begin{definition}
Если функция $f\mathbf 1_D$ интегрируема на $I$, то интеграл по области $D$ определяется равенством
\begin{equation}
\int_D f(x)\,\dd x
:=
\int_I f(x)\mathbf 1_D(x)\,\dd x.
\label{eq:integral-over-domain}
\end{equation}
\end{definition}

Выбор внешнего параллелепипеда $I$ не влияет на результат: вне $D$ характеристическая функция равна нулю. Для областей, ограниченных кусочно-гладкими кривыми или поверхностями, это определение применимо к непрерывным функциям.

\begin{definition}
Объёмом ограниченной допустимой области $D\subset\R^n$ называется число
\begin{equation}
\operatorname{vol}_n(D)
=
\int_D 1\,\dd x.
\label{eq:domain-volume-integral}
\end{equation}
\end{definition}

Формула \eqref{eq:domain-volume-integral} не добавляет искусственную функцию: единица означает постоянную плотность, равную одной единице величины на единицу объёма.

\subsection{Основные свойства}

Пусть функции $f$ и $g$ интегрируемы на допустимой области $D$.

\begin{theorem}[линейность]
Для любых чисел $\alpha,\beta\in\R$
\[
\int_D(\alpha f+\beta g)\,\dd x
=
\alpha\int_Df\,\dd x
+
\beta\int_Dg\,\dd x.
\]
\end{theorem}

Линейность позволяет разбирать сложную плотность на сумму простых вкладов. Например, постоянная и линейная части распределения массы интегрируются отдельно.

\begin{theorem}[аддитивность по области]
Если область $D$ разбита на допустимые части $D_1$ и $D_2$, внутренности которых не пересекаются, то
\[
\int_D f\,\dd x
=
\int_{D_1}f\,\dd x
+
\int_{D_2}f\,\dd x.
\]
\end{theorem}

Общая граница частей не создаёт дополнительного вклада, поскольку её $n$-мерный объём равен нулю.

\begin{theorem}[монотонность и оценки]
Если $f(x)\leq g(x)$ на $D$, то
\[
\int_D f\,\dd x
\leq
\int_D g\,\dd x.
\]
Если
\[
m\leq f(x)\leq M
\quad\text{на }D,
\]
то
\begin{equation}
m\operatorname{vol}_n(D)
\leq
\int_D f\,\dd x
\leq
M\operatorname{vol}_n(D).
\label{eq:multiple-integral-bounds}
\end{equation}
Кроме того,
\[
\abs{\int_Df\,\dd x}
\leq
\int_D\abs{f}\,\dd x.
\]
\end{theorem}

Оценка \eqref{eq:multiple-integral-bounds} служит быстрой проверкой вычисления. Интеграл не может оказаться меньше минимального значения функции, умноженного на объём области, или больше соответствующего максимального значения.

\subsection{Теорема Фубини и повторные интегралы}

Кратный интеграл задаётся единым пределом многомерных сумм. Для вычислений его требуется свести к обычным одномерным интегралам.

\begin{theorem}[Фубини для прямоугольника]
Пусть функция $f$ непрерывна на прямоугольнике
\[
R=[a,b]\times[c,d].
\]
Тогда
\begin{equation}
\iint_R f(x,y)\,\dd x\,\dd y
=
\int_a^b\left(\int_c^d f(x,y)\,\dd y\right)\dd x
=
\int_c^d\left(\int_a^b f(x,y)\,\dd x\right)\dd y.
\label{eq:fubini-rectangle}
\end{equation}
\end{theorem}

В первом повторном интеграле $x$ сначала считается фиксированным. Внутренний интеграл
\[
F(x)=\int_c^d f(x,y)\,\dd y
\]
суммирует значения вдоль вертикального сечения. Затем внешний интеграл суммирует вклады всех сечений по $x$. Во втором порядке используются горизонтальные сечения.

Порядок интегрирования можно выбирать по удобству. Результат один и тот же, но сложность промежуточных вычислений может существенно различаться.

\subsection{Области с переменными пределами}

Рассмотрим область
\begin{equation}
D
=
\{(x,y)\in\R^2
\mid
 a\leq x\leq b,
\ \varphi(x)\leq y\leq\psi(x)\},
\label{eq:x-simple-domain}
\end{equation}
где $\varphi$ и $\psi$ непрерывны и $\varphi(x)\leq\psi(x)$. Такая область называется простой относительно оси $y$ или вертикально простой.

Для фиксированного $x$ вертикальное сечение начинается на нижней границе $y=\varphi(x)$ и заканчивается на верхней границе $y=\psi(x)$. Поэтому
\begin{equation}
\iint_D f(x,y)\,\dd x\,\dd y
=
\int_a^b
\left(
\int_{\varphi(x)}^{\psi(x)}f(x,y)\,\dd y
\right)\dd x.
\label{eq:fubini-x-simple}
\end{equation}

Аналогично, если
\[
D
=
\{(x,y)
\mid
 c\leq y\leq d,
\ \alpha(y)\leq x\leq\beta(y)\},
\]
то
\begin{equation}
\iint_D f(x,y)\,\dd x\,\dd y
=
\int_c^d
\left(
\int_{\alpha(y)}^{\beta(y)}f(x,y)\,\dd x
\right)\dd y.
\label{eq:fubini-y-simple}
\end{equation}

Пределы нельзя выписывать только по рисунку «слева направо». Надёжный порядок таков: сначала выбирается внешний параметр, затем для его фиксированного значения определяется начало и конец внутреннего сечения.

В трёхмерном случае та же идея приводит к записи
\[
\iiint_G f(x,y,z)\,\dd x\,\dd y\,\dd z
=
\int_a^b
\int_{\varphi_1(x)}^{\varphi_2(x)}
\int_{\psi_1(x,y)}^{\psi_2(x,y)}
 f(x,y,z)\,\dd z\,\dd y\,\dd x.
\]
Внутренний интеграл суммирует вклад вдоль отрезка по $z$, средний --- по плоскому сечению, внешний --- по всей области.

\subsection{Физические и статистические интерпретации}

Пусть $D\subset\R^2$ --- плоская пластина, а $\rho(x,y)$ --- поверхностная плотность в $\text{кг}/\text{м}^2$. Тогда масса равна
\begin{equation}
M
=
\iint_D\rho(x,y)\,\dd x\,\dd y.
\label{eq:plate-mass}
\end{equation}
Размерность результата проверяется автоматически:
\[
\frac{\text{кг}}{\text{м}^2}\cdot\text{м}^2
=\text{кг}.
\]

Координаты центра масс определяются формулами
\begin{equation}
\bar x
=
\frac{1}{M}\iint_Dx\rho(x,y)\,\dd x\,\dd y,
\qquad
\bar y
=
\frac{1}{M}\iint_Dy\rho(x,y)\,\dd x\,\dd y.
\label{eq:plate-center-mass}
\end{equation}
В числителе каждой дроби координата умножается на элемент массы $\rho\,\dd x\,\dd y$. Поэтому более тяжёлые участки сильнее смещают центр масс в свою сторону.

Среднее значение скалярного поля $u$ по области равно
\begin{equation}
\langle u\rangle_D
=
\frac{1}{\operatorname{vol}_n(D)}
\int_Du(x)\,\dd x.
\label{eq:spatial-average}
\end{equation}
Если среднее должно учитывать неоднородную массу, используется взвешенная формула
\begin{equation}
\langle u\rangle_{\rho}
=
\frac{1}{M}
\iint_Du(x,y)\rho(x,y)\,\dd x\,\dd y.
\label{eq:mass-weighted-average}
\end{equation}

\begin{problem}[Задача по материалу В.~А.~Зорича: два порядка интегрирования]
На треугольной области
\[
D
=
\{(x,y)\in\R^2
\mid
x\geq0,
\ y\geq0,
\ x+y\leq1\}
\]
вычислить
\[
I=\iint_D(x+2y)\,\dd x\,\dd y
\]
двумя порядками интегрирования и проверить результат оценкой через площадь области.
\end{problem}

\begin{proof}[Решение]
При вертикальных сечениях внешний параметр удовлетворяет $0\leq x\leq1$. Для фиксированного $x$ координата $y$ изменяется от нижней стороны $y=0$ до наклонной стороны $y=1-x$. Поэтому
\[
I
=
\int_0^1\int_0^{1-x}(x+2y)\,\dd y\,\dd x.
\]
Внутренний интеграл равен
\[
\int_0^{1-x}(x+2y)\,\dd y
=
\left[xy+y^2\right]_{0}^{1-x}
=
x(1-x)+(1-x)^2
=1-x.
\]
Следовательно,
\[
I
=
\int_0^1(1-x)\,\dd x
=
\left[x-\frac{x^2}{2}\right]_0^1
=\frac12.
\]

Теперь используем горизонтальные сечения. Имеем $0\leq y\leq1$, а при фиксированном $y$
\[
0\leq x\leq1-y.
\]
Отсюда
\[
I
=
\int_0^1\int_0^{1-y}(x+2y)\,\dd x\,\dd y.
\]
Внутренний интеграл:
\[
\int_0^{1-y}(x+2y)\,\dd x
=
\frac{(1-y)^2}{2}+2y(1-y)
=
\frac12+y-\frac32y^2.
\]
Тогда
\[
I
=
\int_0^1
\left(
\frac12+y-\frac32y^2
\right)\dd y
=
\frac12.
\]

Площадь треугольника равна
\[
\operatorname{area}(D)=\frac12.
\]
На области $D$ функция $x+2y$ принимает значения от $0$ до $2$. Поэтому общая оценка даёт
\[
0
\leq I
\leq
2\cdot\frac12=1.
\]
Полученное значение $I=1/2$ попадает в допустимый диапазон. Кроме того, среднее значение функции по области равно
\[
\frac{I}{\operatorname{area}(D)}=1.
\]
\end{proof}

\begin{problem}[Практическая задача: неоднородная пластина]
Пластина занимает треугольную область
\[
D
=
\left\{
(x,y)
\mid
0\leq x\leq2,
\ 0\leq y\leq1-\frac{x}{2}
\right\},
\]
где координаты измеряются в метрах. Поверхностная плотность и температура заданы функциями
\[
\rho(x,y)=3+x+2y
\quad\frac{\text{кг}}{\text{м}^2},
\]
\[
T(x,y)=20+5x-4y
\quad{}^{\circ}\text{C}.
\]
Нужно найти площадь, массу, центр масс и среднюю по массе температуру пластины.
\end{problem}

\begin{proof}[Решение]
Граница
\[
y=1-\frac{x}{2}
\]
соединяет точки $(0,1)$ и $(2,0)$. При фиксированном $x\in[0,2]$ вертикальное сечение начинается на оси $y=0$ и заканчивается на этой прямой. Поэтому все интегралы имеют пределы
\[
0\leq x\leq2,
\qquad
0\leq y\leq1-\frac{x}{2}.
\]

Площадь равна
\[
A
=
\int_0^2\int_0^{1-x/2}1\,\dd y\,\dd x
=
\int_0^2\left(1-\frac{x}{2}\right)\dd x
=1\ \text{м}^2.
\]
Это совпадает с геометрической формулой для треугольника:
\[
A=\frac12\cdot2\cdot1=1\ \text{м}^2.
\]

Масса вычисляется по формуле \eqref{eq:plate-mass}:
\[
M
=
\int_0^2\int_0^{1-x/2}(3+x+2y)\,\dd y\,\dd x.
\]
Сначала интегрируем по $y$:
\[
\int_0^{1-x/2}(3+x+2y)\,\dd y
=
(3+x)\left(1-\frac{x}{2}\right)
+
\left(1-\frac{x}{2}\right)^2.
\]
После упрощения
\[
(3+x)\left(1-\frac{x}{2}\right)
+
\left(1-\frac{x}{2}\right)^2
=
4-\frac{3x}{2}-\frac{x^2}{4}.
\]
Поэтому
\[
M
=
\int_0^2
\left(
4-\frac{3x}{2}-\frac{x^2}{4}
\right)\dd x
=
\frac{13}{3}\ \text{кг}.
\]

Для координаты $\bar x$ требуется первый момент
\[
M_x^{(1)}
=
\iint_Dx\rho(x,y)\,\dd x\,\dd y.
\]
Здесь верхний индекс напоминает, что это момент первого порядка, а не масса. Подставляем функцию плотности и те же пределы:
\[
M_x^{(1)}
=
\int_0^2
x\left(
4-\frac{3x}{2}-\frac{x^2}{4}
\right)\dd x
=3\ \text{кг}\cdot\text{м}.
\]
Следовательно,
\[
\bar x
=
\frac{M_x^{(1)}}{M}
=
\frac{3}{13/3}
=
\frac{9}{13}\ \text{м}
\approx0.6923\ \text{м}.
\]

Для второй координаты
\[
M_y^{(1)}
=
\int_0^2
\int_0^{1-x/2}
 y(3+x+2y)\,\dd y\,\dd x.
\]
Внутренний интеграл равен
\[
\int_0^{1-x/2}
 y(3+x+2y)\,\dd y
=
\frac{3+x}{2}
\left(1-\frac{x}{2}\right)^2
+
\frac{2}{3}
\left(1-\frac{x}{2}\right)^3.
\]
После внешнего интегрирования получаем
\[
M_y^{(1)}=\frac32\ \text{кг}\cdot\text{м}.
\]
Поэтому
\[
\bar y
=
\frac{M_y^{(1)}}{M}
=
\frac{3/2}{13/3}
=
\frac{9}{26}\ \text{м}
\approx0.3462\ \text{м}.
\]

Средняя по массе температура определяется не обычным средним по площади, а формулой \eqref{eq:mass-weighted-average}:
\[
\overline T_{\rho}
=
\frac{1}{M}
\iint_D T(x,y)\rho(x,y)\,\dd x\,\dd y.
\]
Подставляем обе модели:
\[
Q_T
=
\int_0^2\int_0^{1-x/2}
(20+5x-4y)(3+x+2y)
\,\dd y\,\dd x.
\]
Величина $Q_T$ имеет размерность $^{\circ}\text{C}\cdot\text{кг}$. Последовательное интегрирование даёт
\[
Q_T=\frac{287}{3}\ ^{\circ}\text{C}\cdot\text{кг}.
\]
Следовательно,
\[
\overline T_{\rho}
=
\frac{287/3}{13/3}
=
\frac{287}{13}\ ^{\circ}\text{C}
\approx22.0769\ ^{\circ}\text{C}.
\]

Проверим геометрию результата. Центр масс должен лежать внутри треугольника. Действительно,
\[
0<\frac{9}{13}<2,
\qquad
0<\frac{9}{26}<1-\frac{1}{2}\cdot\frac{9}{13}
=
\frac{17}{26}.
\]
Температура на вершинах равна
\[
T(0,0)=20,
\qquad
T(2,0)=30,
\qquad
T(0,1)=16.
\]
По линейности $T$ её минимум и максимум на треугольнике достигаются в вершинах. Значит,
\[
16\leq\overline T_{\rho}\leq30,
\]
что также выполнено.
\end{proof}


\section{Замена переменных и геометрический смысл якобиана}
\label{sec:change-of-variables}

В предыдущем разделе область интегрирования задавалась непосредственно в декартовых координатах. Такой способ удобен для прямоугольников и областей, ограниченных графиками функций. Для круга, эллипса, кольца или тела вращения декартовы пределы часто разбивают вычисление на несколько частей. Замена переменных позволяет описать геометрию области в более подходящих координатах, но при этом требуется учесть локальное растяжение площади или объёма.

\subsection{Координатное отображение}

Пусть новые переменные собраны в вектор
\[
 t=(t_1,\dots,t_n)\in D_t\subseteq\R^n,
\]
а старые переменные задаются отображением
\[
 x=\Phi(t),
\qquad
\Phi:D_t\to D_x\subseteq\R^n.
\]
В координатах это означает
\[
 x_1=\Phi_1(t_1,\dots,t_n),
\quad\dots\quad,
 x_n=\Phi_n(t_1,\dots,t_n).
\]
Область $D_t$ выбирают так, чтобы её границы были проще границ исходной области $D_x$.

\begin{example}
Эллипс
\[
D_x=
\left\{
(x,y)
\mid
\frac{x^2}{a^2}+\frac{y^2}{b^2}\leq1
\right\}
\]
получается из единичного круга
\[
D_t=
\{(u,v)\mid u^2+v^2\leq1\}
\]
при линейной замене
\[
x=au,
\qquad
y=bv.
\]
В новых координатах сложная граница эллипса превращается в окружность.
\end{example}

Для корректной замены недостаточно просто переписать функцию через новые переменные. Малый квадрат в координатах $t$ после отображения может превратиться в параллелограмм другой площади. Именно этот эффект описывает определитель якобиана.

\subsection{Линейное растяжение площади и объёма}

Сначала рассмотрим линейное отображение
\[
x=Lt,
\qquad
L\in\R^{n\times n}.
\]
Столбец $Le_j$ показывает, во что переходит единичный отрезок, направленный вдоль $j$-й координатной оси. Поэтому столбцы матрицы $L$ являются рёбрами образа единичного куба.

\begin{theorem}[масштабирование объёма линейным отображением]
Если $E\subseteq\R^n$ является допустимым измеримым множеством, то
\[
\operatorname{vol}(L(E))
=
\abs{\det L}\operatorname{vol}(E).
\]
\end{theorem}

В двумерном случае $\abs{\det L}$ равен площади параллелограмма, построенного на столбцах матрицы $L$. В трёхмерном случае это объём параллелепипеда. Знак определителя отвечает за ориентацию, но площадь и объём не могут быть отрицательными, поэтому берётся модуль.

\begin{example}
Для преобразования
\[
\begin{pmatrix}x\\y\end{pmatrix}
=
\begin{pmatrix}
 a&0\\
 0&b
\end{pmatrix}
\begin{pmatrix}u\\v\end{pmatrix}
\]
имеем
\[
\det L=ab.
\]
Следовательно, площадь единичного круга, равная $\pi$, после растяжения по осям в $a$ и $b$ раз становится
\[
\pi\abs{ab}.
\]
При $a,b>0$ получаем обычную формулу площади эллипса $\pi ab$.
\end{example}

\subsection{Якобиан как локальный коэффициент растяжения}

Пусть отображение $\Phi$ дифференцируемо в точке $t$. Для малого приращения $h$ справедливо приближение
\[
\Phi(t+h)-\Phi(t)
=
\Jac{\Phi}(t)h+o(\norm{h}).
\]
На достаточно малом участке нелинейное отображение ведёт себя почти как линейное отображение с матрицей $\Jac{\Phi}(t)$. Поэтому локальный коэффициент изменения $n$-мерного объёма равен
\[
\abs{\det\Jac{\Phi}(t)}.
\]

\begin{definition}
Определителем якобиана замены $x=\Phi(t)$ называется функция
\[
J_{\Phi}(t)
=
\det\Jac{\Phi}(t)
=
\det
\begin{pmatrix}
\pdv{\Phi_1}{t_1}&\cdots&\pdv{\Phi_1}{t_n}\\
\vdots&\ddots&\vdots\\
\pdv{\Phi_n}{t_1}&\cdots&\pdv{\Phi_n}{t_n}
\end{pmatrix}.
\]
\end{definition}

Строка $i$ содержит производные старой координаты $x_i=\Phi_i(t)$ по всем новым переменным. Столбец $j$ показывает, как меняется весь вектор $x$ при изменении только $t_j$. Определитель объединяет эти направленные изменения в один коэффициент растяжения объёма.

\begin{caution}
Матрица Якоби и её определитель выполняют разные роли. Матрица описывает локальное линейное преобразование всех направлений. В формулу замены переменных в кратном интеграле входит модуль её определителя.
\end{caution}

\subsection{Формула замены переменных}

\begin{theorem}[замена переменных в кратном интеграле]
Пусть $\Phi:D_t\to D_x$ является взаимно однозначным отображением класса $C^1$ во внутренних точках областей, причём
\[
\det\Jac{\Phi}(t)\neq0.
\]
Пусть возможные нарушения взаимной однозначности или гладкости находятся только на границах либо на множествах нулевого объёма. Тогда для интегрируемой функции $f$ справедливо
\begin{equation}
\label{eq:change-of-variables}
\int_{D_x}f(x)\,\dd x
=
\int_{D_t}
 f(\Phi(t))
 \abs{\det\Jac{\Phi}(t)}
\,\dd t.
\end{equation}
\end{theorem}

В двумерном случае формула принимает вид
\[
\iint_{D_x}f(x,y)\,\dd x\,\dd y
=
\iint_{D_t}
 f\bigl(x(u,v),y(u,v)\bigr)
 \abs{\pdv{(x,y)}{(u,v)}}
\,\dd u\,\dd v,
\]
где
\[
\pdv{(x,y)}{(u,v)}
=
\det
\begin{pmatrix}
\pdv{x}{u}&\pdv{x}{v}\\
\pdv{y}{u}&\pdv{y}{v}
\end{pmatrix}.
\]

Формула состоит из трёх независимых замен:
\[
f(x)\longmapsto f(\Phi(t)),
\]
\[
D_x\longmapsto D_t,
\]
\[
\dd x\longmapsto
\abs{\det\Jac{\Phi}(t)}\,\dd t.
\]
Пропуск любого из этих трёх шагов приводит к неверному интегралу.

\begin{proofidea}
Разобьём $D_t$ на малые ячейки. В окрестности точки $t$ отображение $\Phi$ почти совпадает со своей линейной частью $\Jac{\Phi}(t)$. Поэтому объём образа малой ячейки приближённо равен её объёму, умноженному на $\abs{\det\Jac{\Phi}(t)}$. После подстановки $f(\Phi(t))$ получаются интегральные суммы для правой части формулы \eqref{eq:change-of-variables}. Предельный переход даёт равенство интегралов.
\end{proofidea}

В вычислении сначала задают отображение \(x=\Phi(t)\) и новую область, затем находят \(\Jac{\Phi}(t)\) и \(\abs{\det\Jac{\Phi}(t)}\), подставляют \(x=\Phi(t)\) в подынтегральную функцию и лишь после записи новых пределов вычисляют интеграл.

\subsection{Полярные координаты}

Полярные координаты на плоскости задаются формулами
\begin{equation}
\label{eq:polar-coordinates}
x=r\cos\varphi,
\qquad
y=r\sin\varphi.
\end{equation}
Здесь $r\geq0$ --- расстояние до начала координат, а $\varphi$ --- угол.

Матрица Якоби равна
\[
\Jac{\Phi}(r,\varphi)
=
\begin{pmatrix}
\cos\varphi&-r\sin\varphi\\
\sin\varphi&r\cos\varphi
\end{pmatrix}.
\]
Её определитель
\[
J_{\Phi}(r,\varphi)
=
r\cos^2\varphi+r\sin^2\varphi
=r.
\]
Следовательно,
\begin{equation}
\label{eq:polar-area-element}
\dd x\,\dd y
=
r\,\dd r\,\dd\varphi.
\end{equation}

Множитель $r$ имеет геометрический смысл. Малое изменение угла $\Delta\varphi$ на расстоянии $r$ создаёт дугу длины примерно $r\Delta\varphi$. Вместе с радиальной толщиной $\Delta r$ получается площадь
\[
\Delta A
\approx
\Delta r\cdot r\Delta\varphi.
\]

В точке $r=0$ якобиан обращается в нуль, а угловая координата не определяется однозначно. Это происходит на множестве нулевой площади и не мешает использованию формулы интегрирования по кругу.

\begin{example}
Для круга
\[
K_R=
\{(x,y)\mid x^2+y^2\leq R^2\}
\]
новая область имеет прямоугольное описание
\[
0\leq r\leq R,
\qquad
0\leq\varphi\leq2\pi.
\]
Поэтому
\[
\iint_{K_R}f(x,y)\,\dd x\,\dd y
=
\int_0^{2\pi}\int_0^R
f(r\cos\varphi,r\sin\varphi)
 r\,\dd r\,\dd\varphi.
\]
\end{example}

\subsection{Цилиндрические и сферические координаты}

Цилиндрические координаты сохраняют высоту $z$ и используют полярные координаты в плоскости:
\[
x=r\cos\varphi,
\qquad
y=r\sin\varphi,
\qquad z=z.
\]
Определитель якобиана равен $r$, поэтому
\[
\dd x\,\dd y\,\dd z
=
r\,\dd r\,\dd\varphi\,\dd z.
\]
Эта система удобна для цилиндров, тел вращения и областей с осевой симметрией.

В сферических координатах примем соглашение
\[
x=\rho\sin\theta\cos\varphi,
\]
\[
y=\rho\sin\theta\sin\varphi,
\]
\[
z=\rho\cos\theta,
\]
где
\[
\rho\geq0,
\qquad
0\leq\theta\leq\pi,
\qquad
0\leq\varphi<2\pi.
\]
Здесь $\theta$ отсчитывается от положительного направления оси $z$. Определитель якобиана по модулю равен
\[
\rho^2\sin\theta,
\]
поэтому
\begin{equation}
\label{eq:spherical-volume-element}
\dd x\,\dd y\,\dd z
=
\rho^2\sin\theta
\,\dd\rho\,\dd\theta\,\dd\varphi.
\end{equation}

Два множителя имеют разный смысл: $\rho^2$ отражает квадратичный рост площади сферического слоя, а $\sin\theta$ учитывает сжатие окружностей широты около полюсов.

\begin{problem}[Задача по материалу В.~А.~Зорича]
С помощью полярных координат вычислить интеграл
\[
I
=
\iint_{x^2+y^2\leq R^2}
(x^2+y^2)\,\dd x\,\dd y,
\qquad R>0.
\]
\end{problem}

\begin{proof}[Решение]
Исходная область является кругом радиуса $R$, поэтому используем замену \eqref{eq:polar-coordinates}. Подынтегральная функция упрощается:
\[
x^2+y^2
=
r^2\cos^2\varphi+r^2\sin^2\varphi
=r^2.
\]
Элемент площади по формуле \eqref{eq:polar-area-element} равен
\[
\dd x\,\dd y=r\,\dd r\,\dd\varphi.
\]
Следовательно,
\[
I
=
\int_0^{2\pi}\int_0^R
r^2\cdot r\,\dd r\,\dd\varphi.
\]
Важно, что третья степень $r$ появилась из двух разных источников: $r^2$ пришло из функции, а дополнительный множитель $r$ --- из якобиана.

Вычисляем внутренний интеграл:
\[
\int_0^Rr^3\,\dd r
=
\frac{R^4}{4}.
\]
Поэтому
\[
I
=
\int_0^{2\pi}\frac{R^4}{4}\,\dd\varphi
=
\frac{\pi R^4}{2}.
\]

Проверим размерность. Величина $x^2+y^2$ имеет размерность длины в квадрате, а элемент площади --- также длины в квадрате. Поэтому интеграл должен иметь размерность длины в четвёртой степени, что согласуется с $R^4$.
\end{proof}

\begin{problem}[Практическая задача: эллиптическое пятно нагрева]
Источник тепла действует на эллиптическую область
\[
D=
\left\{
(x,y)
\mid
\frac{x^2}{a^2}+\frac{y^2}{b^2}\leq1
\right\},
\qquad a,b>0,
\]
где $a$ и $b$ измеряются в метрах. Поверхностная плотность мощности равна
\[
q(x,y)
=
q_c
\left(
1-\frac{x^2}{a^2}-\frac{y^2}{b^2}
\right),
\qquad
q_c>0,
\]
и измеряется в ваттах на квадратный метр.

Нужно найти полную мощность источника, среднюю плотность мощности по области и размеры соосного эллипса, внутри которого выделяется доля $\eta\in(0,1)$ полной мощности.
\end{problem}

\begin{proof}[Решение]
Граница области содержит выражение
\[
\frac{x^2}{a^2}+\frac{y^2}{b^2}.
\]
Поэтому вводим эллиптические полярные координаты
\begin{equation}
\label{eq:elliptic-polar}
x=ar\cos\varphi,
\qquad
y=br\sin\varphi.
\end{equation}
Тогда
\[
\frac{x^2}{a^2}+\frac{y^2}{b^2}
=
r^2,
\]
а исходный эллипс переходит в прямоугольник параметров
\[
0\leq r\leq1,
\qquad
0\leq\varphi\leq2\pi.
\]

Матрица Якоби равна
\[
\Jac{\Phi}(r,\varphi)
=
\begin{pmatrix}
 a\cos\varphi&-ar\sin\varphi\\
 b\sin\varphi&br\cos\varphi
\end{pmatrix}.
\]
Её определитель
\[
\det\Jac{\Phi}
=
abr\cos^2\varphi
+abr\sin^2\varphi
=abr.
\]
Так как $a,b,r\geq0$, модуль определителя также равен $abr$. Размерность $ab$ равна $\text{м}^2$, а $r$ безразмерен, поэтому якобиан имеет правильную размерность площади.

После замены функция мощности становится
\[
q(\Phi(r,\varphi))
=q_c(1-r^2).
\]
Полная мощность равна
\[
P
=
\iint_Dq(x,y)\,\dd x\,\dd y
=
\int_0^{2\pi}\int_0^1
q_c(1-r^2)abr
\,\dd r\,\dd\varphi.
\]
Выносим постоянные множители:
\[
P
=2\pi abq_c
\int_0^1(r-r^3)\,\dd r.
\]
Внутренний интеграл равен
\[
\int_0^1(r-r^3)\,\dd r
=
\frac12-\frac14
=
\frac14.
\]
Следовательно,
\begin{equation}
\label{eq:total-elliptic-power}
P
=
\frac{\pi abq_c}{2}.
\end{equation}
Размерность результата --- ватт, поскольку
\[
\text{м}^2\cdot\frac{\text{Вт}}{\text{м}^2}
=
\text{Вт}.
\]

Площадь эллипса можно получить той же заменой при подынтегральной функции $1$:
\[
A
=
\int_0^{2\pi}\int_0^1abr\,\dd r\,\dd\varphi
=
\pi ab.
\]
Поэтому средняя по площади плотность мощности равна
\[
\overline q
=
\frac{P}{A}
=
\frac{\pi abq_c/2}{\pi ab}
=
\frac{q_c}{2}.
\]

Рассмотрим соосный эллипс, соответствующий условию
\[
0\leq r\leq s,
\qquad 0<s<1.
\]
Его полуоси равны $as$ и $bs$. Мощность внутри него составляет
\[
P(s)
=
2\pi abq_c
\int_0^s(r-r^3)\,\dd r.
\]
Следовательно,
\[
P(s)
=
2\pi abq_c
\left(
\frac{s^2}{2}-\frac{s^4}{4}
\right)
=
\pi abq_c
\left(
s^2-\frac{s^4}{2}
\right).
\]
Доля полной мощности равна
\[
\frac{P(s)}{P}
=
\frac{
\pi abq_c(s^2-s^4/2)
}{
\pi abq_c/2
}
=
2s^2-s^4.
\]
Требование накопить долю $\eta$ приводит к уравнению
\[
2s^2-s^4=\eta.
\]
Положим $z=s^2$. Тогда
\[
z^2-2z+\eta=0,
\]
откуда
\[
z=1\pm\sqrt{1-\eta}.
\]
Поскольку $0<z<1$, выбираем знак минус:
\[
s^2=1-\sqrt{1-\eta},
\]
\begin{equation}
\label{eq:power-fraction-radius}
s
=
\sqrt{1-\sqrt{1-\eta}}.
\end{equation}
Искомый эллипс имеет полуоси
\[
a_{\eta}
=
a\sqrt{1-\sqrt{1-\eta}},
\]
\[
b_{\eta}
=
b\sqrt{1-\sqrt{1-\eta}}.
\]

Например, для $\eta=0.9$ получаем
\[
s
=
\sqrt{1-\sqrt{0.1}}
\approx0.8269.
\]
Значит, $90\%$ мощности выделяется внутри эллипса с полуосями примерно $0.827a$ и $0.827b$. Это меньше, чем $90\%$ линейных размеров, потому что плотность мощности максимальна в центре и убывает к границе.
\end{proof}


\section{Векторный анализ}
\label{sec:vector-analysis}

Векторный анализ связывает локальные производные поля с интегральными характеристиками кривых, поверхностей и пространственных областей. Его основные объекты --- скалярное поле
\[
f\colon D\subseteq\R^3\to\R
\]
и векторное поле
\[
F\colon D\subseteq\R^3\to\R^3,
\qquad
F=(P,Q,R).
\]
Скалярное поле может задавать температуру, давление или концентрацию, а векторное --- скорость потока, силу или плотность теплового потока.

\subsection{Градиент, дивергенция и ротор}

Для гладкого скалярного поля градиент уже был определён как
\[
\grad f
=
\left(
\frac{\partial f}{\partial x},
\frac{\partial f}{\partial y},
\frac{\partial f}{\partial z}
\right).
\]
Он указывает направление наиболее быстрого роста функции, а его скалярное произведение с единичным вектором $e$ даёт производную по направлению:
\[
D_e f=\grad f\cdot e.
\]

\begin{definition}
Дивергенцией гладкого векторного поля $F=(P,Q,R)$ называется скалярное поле
\begin{equation}
\label{eq:divergence-definition}
\operatorname{div}F
=
\frac{\partial P}{\partial x}
+
\frac{\partial Q}{\partial y}
+
\frac{\partial R}{\partial z}.
\end{equation}
\end{definition}

Дивергенция измеряет локальную интенсивность источника. Если $F$ --- поле скорости с размерностью $\text{м}/\text{с}$, то каждая производная в \eqref{eq:divergence-definition} имеет размерность $1/\text{с}$. Положительная дивергенция означает локальное расширение потока, отрицательная --- сток, а нулевая соответствует отсутствию объёмного источника.

\begin{definition}
Ротором гладкого векторного поля $F=(P,Q,R)$ называется векторное поле
\begin{equation}
\label{eq:curl-definition}
\operatorname{rot}F
=
\left(
\frac{\partial R}{\partial y}-\frac{\partial Q}{\partial z},
\frac{\partial P}{\partial z}-\frac{\partial R}{\partial x},
\frac{\partial Q}{\partial x}-\frac{\partial P}{\partial y}
\right).
\end{equation}
\end{definition}

Ротор измеряет локальную циркуляцию поля. Для поля скорости он имеет размерность $1/\text{с}$ и характеризует ось и интенсивность малого вращения среды. Направление ротора определяется правилом правого винта.

Операторы удовлетворяют двум важным тождествам:
\begin{equation}
\label{eq:curl-grad-zero}
\operatorname{rot}(\grad f)=0,
\end{equation}
\begin{equation}
\label{eq:div-curl-zero}
\operatorname{div}(\operatorname{rot}F)=0.
\end{equation}
Первое равенство следует из равенства смешанных производных. Второе проверяется тем же способом: каждая производная второго порядка входит дважды с противоположными знаками.

\subsection{Криволинейный интеграл и работа поля}

Пусть ориентированная гладкая кривая задана параметризацией
\[
r(t)=(x(t),y(t),z(t)),
\qquad
a\leq t\leq b.
\]
Вектор $r'(t)$ направлен по движению вдоль кривой, а малое перемещение записывается как
\[
\dd r=r'(t)\,\dd t.
\]

\begin{definition}
Работой поля $F$ вдоль ориентированной кривой $\gamma$ называется криволинейный интеграл
\begin{equation}
\label{eq:line-integral-work}
\int_{\gamma}F\cdot\dd r
=
\int_a^b
F(r(t))\cdot r'(t)\,\dd t.
\end{equation}
\end{definition}

Формула \eqref{eq:line-integral-work} задаёт вычислительный алгоритм. Сначала подставляют координаты пути в компоненты поля, затем находят производную $r'(t)$, вычисляют скалярное произведение и только после этого интегрируют по параметру.

Если $F$ имеет размерность силы, а $r$ --- длины, то интеграл имеет размерность работы. При замене параметра значение интеграла не меняется, если сохраняются сама ориентированная кривая и направление её обхода.

\begin{definition}
Поле $F$ называется потенциальным в области $D$, если существует функция $U\colon D\to\R$, для которой
\[
F=\grad U.
\]
Функция $U$ называется потенциалом поля.
\end{definition}

Для потенциального поля многомерная формула Ньютона--Лейбница принимает вид
\begin{equation}
\label{eq:fundamental-line-integral}
\int_{\gamma}\grad U\cdot\dd r
=
U(r(b))-U(r(a)).
\end{equation}
Следовательно, работа зависит только от начальной и конечной точек. Циркуляция потенциального поля по любому замкнутому пути равна нулю.

Обратно, непрерывное поле потенциально тогда и только тогда, когда его интеграл по каждому замкнутому пути равен нулю. Для гладкого поля в выпуклой области достаточно проверить условие
\[
\operatorname{rot}F=0.
\]
Оговорка об области существенна: отсутствие локального вихря не устраняет глобальные препятствия, например выколотую точку, вокруг которой можно совершить ненулевой обход.

\subsection{Поток через поверхность}

Пусть ориентированная поверхность задана параметризацией
\[
r(u,v),
\qquad
(u,v)\in G\subseteq\R^2.
\]
Векторы $r_u$ и $r_v$ касаются поверхности. Их векторное произведение
\[
r_u\times r_v
\]
перпендикулярно поверхности и задаёт выбранную ориентацию. Его длина является коэффициентом перехода от площади в плоскости параметров к площади на поверхности.

\begin{definition}
Потоком поля $F$ через ориентированную поверхность $S$ называется интеграл
\begin{equation}
\label{eq:surface-flux-parametric}
\iint_S F\cdot n\,\dd\sigma
=
\iint_G
F(r(u,v))\cdot
\bigl(r_u\times r_v\bigr)
\,\dd u\,\dd v,
\end{equation}
где $n$ --- единичная нормаль, согласованная с ориентацией поверхности.
\end{definition}

В правой части отдельно нормировать вектор $r_u\times r_v$ не требуется: его направление задаёт нормаль, а длина уже содержит коэффициент изменения площади. Смена ориентации меняет знак потока.

Если $F$ --- скорость жидкости, то поток через поверхность имеет размерность
\[
\frac{\text{м}}{\text{с}}\cdot\text{м}^2
=
\frac{\text{м}^3}{\text{с}},
\]
то есть представляет объёмный расход.

\subsection{Формулы Грина, Стокса и Гаусса}

Три классические формулы связывают интегралы по границе с производными поля внутри области.

Для положительно ориентированной границы $\partial D$ плоской области $D$ формула Грина имеет вид
\begin{equation}
\label{eq:green-theorem}
\oint_{\partial D}
P\,\dd x+Q\,\dd y
=
\iint_D
\left(
\frac{\partial Q}{\partial x}
-
\frac{\partial P}{\partial y}
\right)
\dd x\,\dd y.
\end{equation}
Положительная ориентация означает обход против часовой стрелки, при котором область остаётся слева.

Для ориентированной поверхности $S$ с согласованно ориентированной границей $\partial S$ формула Стокса записывается как
\begin{equation}
\label{eq:stokes-vector}
\oint_{\partial S}F\cdot\dd r
=
\iint_S
\operatorname{rot}F\cdot n\,\dd\sigma.
\end{equation}
Она утверждает: циркуляция поля по краю равна потоку ротора через поверхность.

Для пространственной области $\Omega$ с внешней ориентацией границы $\partial\Omega$ формула Гаусса--Остроградского имеет вид
\begin{equation}
\label{eq:gauss-divergence}
\iint_{\partial\Omega}
F\cdot n\,\dd\sigma
=
\iiint_{\Omega}
\operatorname{div}F\,\dd V.
\end{equation}
Она утверждает: полный поток наружу через замкнутую поверхность равен суммарной мощности источников внутри области.

Формула Грина является плоским вариантом формулы Стокса. Все три равенства имеют общую структуру: интеграл производной по области равен интегралу исходного объекта по её ориентированной границе.

\begin{problem}[Задача по материалу Зорича: поле жёсткого вращения]
Тело вращается вокруг оси $Oz$ с постоянной угловой скоростью $\omega$. Поле скоростей равно
\[
v(x,y,z)=(-\omega y,\omega x,0).
\]
Найти дивергенцию и ротор поля, вычислить циркуляцию по окружности
\[
x^2+y^2=R^2,
\qquad z=0,
\]
ориентированной против часовой стрелки, и проверить формулу Стокса.
\end{problem}

\begin{proof}[Решение]
Обозначим компоненты поля через
\[
P=-\omega y,
\qquad
Q=\omega x,
\qquad
R_3=0.
\]
Индекс у $R_3$ нужен, чтобы не спутать третью компоненту поля с радиусом окружности.

По определению дивергенции
\[
\operatorname{div}v
=
\frac{\partial(-\omega y)}{\partial x}
+
\frac{\partial(\omega x)}{\partial y}
+
\frac{\partial0}{\partial z}
=0.
\]
Следовательно, жёсткое вращение не создаёт локального объёмного расширения.

Для ротора вычисляем компоненты по формуле \eqref{eq:curl-definition}:
\[
\operatorname{rot}v
=
\left(
0-0,
0-0,
\frac{\partial(\omega x)}{\partial x}
-
\frac{\partial(-\omega y)}{\partial y}
\right).
\]
Так как
\[
\frac{\partial(\omega x)}{\partial x}=\omega,
\qquad
\frac{\partial(-\omega y)}{\partial y}=-\omega,
\]
получаем
\begin{equation}
\label{eq:rigid-rotation-curl}
\operatorname{rot}v=(0,0,2\omega).
\end{equation}
Ротор направлен вдоль оси вращения, а его модуль равен удвоенной угловой скорости.

Параметризуем окружность:
\[
r(t)=(R\cos t,R\sin t,0),
\qquad
0\leq t\leq2\pi.
\]
Тогда
\[
r'(t)=(-R\sin t,R\cos t,0),
\]
а поле на окружности равно
\[
v(r(t))
=
(-\omega R\sin t,\omega R\cos t,0).
\]
Скалярное произведение принимает вид
\[
v(r(t))\cdot r'(t)
=
\omega R^2\sin^2t+
\omega R^2\cos^2t
=
\omega R^2.
\]
Поэтому циркуляция равна
\begin{equation}
\label{eq:rigid-rotation-circulation}
\oint_{\partial S}v\cdot\dd r
=
\int_0^{2\pi}\omega R^2\,\dd t
=
2\pi\omega R^2.
\end{equation}

Проверим правую часть формулы Стокса. В качестве $S$ берём круг радиуса $R$ в плоскости $z=0$ с нормалью $n=(0,0,1)$. Из \eqref{eq:rigid-rotation-curl}
\[
\operatorname{rot}v\cdot n=2\omega.
\]
Следовательно,
\[
\iint_S
\operatorname{rot}v\cdot n\,\dd\sigma
=
2\omega\,\operatorname{area}(S)
=
2\omega\pi R^2.
\]
Это совпадает с \eqref{eq:rigid-rotation-circulation}.

Если $\omega$ измеряется в $1/\text{с}$, то циркуляция имеет размерность
\[
\frac1{\text{с}}\cdot\text{м}^2
=
\frac{\text{м}^2}{\text{с}},
\]
что согласуется с произведением скорости на длину.
\end{proof}

\begin{problem}[Практическая задача: баланс воздуха в камере]
Прямоугольная камера занимает область
\[
\Omega=
\left[-\frac{L_x}{2},\frac{L_x}{2}\right]
\times
\left[-\frac{L_y}{2},\frac{L_y}{2}\right]
\times
\left[-\frac{L_z}{2},\frac{L_z}{2}\right].
\]
Поле скорости воздуха моделируется формулой
\[
u(x,y,z)=(\alpha x,\beta y,\gamma z),
\]
где $\alpha,\beta,\gamma$ измеряются в $1/\text{с}$.

Для значений
\[
L_x=6\ \text{м},
\qquad
L_y=4\ \text{м},
\qquad
L_z=3\ \text{м},
\]
\[
\alpha=0.020\ \text{с}^{-1},
\qquad
\beta=-0.010\ \text{с}^{-1},
\qquad
\gamma=0.005\ \text{с}^{-1}
\]
нужно найти чистый объёмный расход через границу камеры, проверить ответ прямым суммированием потоков через грани и определить, как изменить $\gamma$, чтобы модель стала несжимаемой.
\end{problem}

\begin{proof}[Решение]
Сначала вычислим локальную дивергенцию:
\[
\operatorname{div}u
=
\frac{\partial(\alpha x)}{\partial x}
+
\frac{\partial(\beta y)}{\partial y}
+
\frac{\partial(\gamma z)}{\partial z}
=
\alpha+\beta+\gamma.
\]
Подставляя параметры, получаем
\[
\operatorname{div}u
=0.020-0.010+0.005
=0.015\ \text{с}^{-1}.
\]
Положительный знак означает, что модель содержит объёмное расширение воздуха.

Объём камеры равен
\[
\abs{\Omega}
=L_xL_yL_z
=6\cdot4\cdot3
=72\ \text{м}^3.
\]
По формуле Гаусса--Остроградского чистый поток наружу равен
\[
Q_{\mathrm{out}}
=
\iiint_{\Omega}
\operatorname{div}u\,\dd V.
\]
Дивергенция постоянна, поэтому
\begin{equation}
\label{eq:chamber-net-flow}
Q_{\mathrm{out}}
=(\alpha+\beta+\gamma)L_xL_yL_z
=0.015\cdot72
=1.08\ \frac{\text{м}^3}{\text{с}}.
\end{equation}
Размерность проверяется непосредственно:
\[
\frac1{\text{с}}\cdot\text{м}^3
=
\frac{\text{м}^3}{\text{с}}.
\]

Теперь вычислим поток по граням. На двух гранях, перпендикулярных оси $Ox$, внешние нормали противоположны, но и значения компоненты $u_x=\alpha x$ имеют противоположные знаки. Поэтому оба вклада складываются:
\[
Q_x
=2\left(\alpha\frac{L_x}{2}\right)L_yL_z
=\alpha L_xL_yL_z
=0.020\cdot72
=1.44\ \frac{\text{м}^3}{\text{с}}.
\]
Аналогично,
\[
Q_y
=\beta L_xL_yL_z
=-0.010\cdot72
=-0.72\ \frac{\text{м}^3}{\text{с}},
\]
\[
Q_z
=\gamma L_xL_yL_z
=0.005\cdot72
=0.36\ \frac{\text{м}^3}{\text{с}}.
\]
Отрицательное значение $Q_y$ означает чистый приток через пару граней, перпендикулярных оси $Oy$. Суммарный поток
\[
Q_x+Q_y+Q_z
=1.44-0.72+0.36
=1.08\ \frac{\text{м}^3}{\text{с}},
\]
что совпадает с \eqref{eq:chamber-net-flow}.

Для несжимаемой стационарной модели требуется
\[
\operatorname{div}u=0.
\]
Следовательно, новое значение коэффициента $\gamma$ должно удовлетворять
\[
\alpha+\beta+\gamma_{\mathrm{new}}=0.
\]
Отсюда
\[
\gamma_{\mathrm{new}}
=-(\alpha+\beta)
=-(0.020-0.010)
=-0.010\ \text{с}^{-1}.
\]
Необходимое изменение равно
\[
\Delta\gamma
=\gamma_{\mathrm{new}}-\gamma
=-0.010-0.005
=-0.015\ \text{с}^{-1}.
\]
После такой коррекции положительный поток через $x$-грани будет точно компенсирован притоком через $y$- и $z$-грани.
\end{proof}


\section{Дифференциальные формы: единый язык анализа}
\label{sec:differential-forms-bridge}

В предыдущем разделе градиент, ротор и дивергенция вводились как отдельные операции над полями, а работа, циркуляция и поток --- как разные типы интегралов. Дифференциальные формы показывают, что это проявления одной конструкции. Здесь нужен лишь рабочий мост к этому языку; полная теория форм и многообразий выходит за рамки главы.

\subsection{Формы степени от нуля до трёх}

Гладкую функцию
\[
f\colon D\subseteq\R^3\to\R
\]
называют дифференциальной формой степени ноль, или кратко \(0\)-формой. Её обычный дифференциал
\[
\dd f
=
\frac{\partial f}{\partial x}\,\dd x
+
\frac{\partial f}{\partial y}\,\dd y
+
\frac{\partial f}{\partial z}\,\dd z
\]
является \(1\)-формой. В каждой точке он принимает вектор малого смещения
\(h=(h_x,h_y,h_z)\) и возвращает линейную часть изменения функции:
\[
Df(h)
=
\frac{\partial f}{\partial x}h_x
+
\frac{\partial f}{\partial y}h_y
+
\frac{\partial f}{\partial z}h_z.
\]

Общая \(1\)-форма в декартовых координатах имеет вид
\begin{equation}
\label{eq:one-form-coordinate}
\omega
=P\,\dd x+Q\,\dd y+R\,\dd z.
\end{equation}
Если \(F=(P,Q,R)\), то
\[
\omega(h)=F\cdot h.
\]
Поэтому \eqref{eq:one-form-coordinate} является формой элементарной работы поля \(F\): её интеграл по ориентированной кривой равен криволинейному интегралу \(\int_C F\cdot\dd r\).

Чтобы описывать ориентированные площадки, вводят внешнее произведение. Для координатных дифференциалов выполняются правила
\[
\dd x\wedge\dd y=-\dd y\wedge\dd x,
\qquad
\dd x\wedge\dd x=0,
\]
и аналогичные соотношения для остальных координат. Перестановка двух направлений меняет ориентацию площадки, поэтому меняется знак формы; совпадающие направления не натягивают площадь, поэтому результат равен нулю.

Общая \(2\)-форма в \(\R^3\) записывается как
\begin{equation}
\label{eq:two-form-coordinate}
\eta
=P\,\dd y\wedge\dd z
+Q\,\dd z\wedge\dd x
+R\,\dd x\wedge\dd y.
\end{equation}
На двух векторах \(u,v\) она принимает значение
\[
\eta(u,v)=F\cdot(u\times v),
\qquad
F=(P,Q,R).
\]
Следовательно, \(2\)-форма естественно описывает элементарный поток через ориентированную площадку.

Наконец, \(3\)-форма имеет вид
\[
\mu=\rho(x,y,z)\,\dd x\wedge\dd y\wedge\dd z.
\]
Она описывает ориентированную объёмную плотность. Её интеграл по области совпадает с обычным тройным интегралом от \(\rho\).

\subsection{Внешний дифференциал}

Операция внешнего дифференцирования обозначается той же буквой \(\dd\), что и обычный дифференциал функции. Она повышает степень формы на единицу:
\[
0\text{-форма}
\xrightarrow{\ \dd\ }
1\text{-форма}
\xrightarrow{\ \dd\ }
2\text{-форма}
\xrightarrow{\ \dd\ }
3\text{-форма}.
\]
Для \(1\)-формы \(\omega=P\,\dd x+Q\,\dd y+R\,\dd z\) непосредственное вычисление даёт
\begin{equation}
\label{eq:exterior-derivative-one-form}
\begin{aligned}
\dd\omega
={}&
\left(
\frac{\partial R}{\partial y}
-
\frac{\partial Q}{\partial z}
\right)\dd y\wedge\dd z
\\
&+
\left(
\frac{\partial P}{\partial z}
-
\frac{\partial R}{\partial x}
\right)\dd z\wedge\dd x
\\
&+
\left(
\frac{\partial Q}{\partial x}
-
\frac{\partial P}{\partial y}
\right)\dd x\wedge\dd y.
\end{aligned}
\end{equation}
Коэффициенты справа являются компонентами \(\operatorname{rot}F\). Значит, внешний дифференциал переводит форму работы поля в форму потока его ротора.

Для \(2\)-формы \eqref{eq:two-form-coordinate} получаем
\begin{equation}
\label{eq:exterior-derivative-two-form}
\dd\eta
=
\left(
\frac{\partial P}{\partial x}
+
\frac{\partial Q}{\partial y}
+
\frac{\partial R}{\partial z}
\right)
\dd x\wedge\dd y\wedge\dd z.
\end{equation}
Коэффициент в скобках равен \(\operatorname{div}F\). Таким образом, одна операция \(\dd\) скрывает знакомую цепочку
\[
f
\longmapsto
\operatorname{grad}f
\longmapsto
\operatorname{rot}F
\longmapsto
\operatorname{div}F.
\]

Если коэффициенты дважды непрерывно дифференцируемы, то
\begin{equation}
\label{eq:d-squared-zero}
\dd(\dd\omega)=0.
\end{equation}
Причина состоит в равенстве смешанных производных и антисимметрии внешнего произведения. Для векторного анализа из \eqref{eq:d-squared-zero} сразу следуют уже встречавшиеся тождества
\[
\operatorname{rot}(\grad f)=0,
\qquad
\operatorname{div}(\operatorname{rot}F)=0.
\]

\subsection{Общая формула Стокса}

Пусть \(S\) --- ориентированная \(k\)-мерная поверхность с краем \(\partial S\), а \(\omega\) --- гладкая форма степени \(k-1\). При согласованной ориентации поверхности и её края выполняется общая формула Стокса
\begin{equation}
\label{eq:general-stokes}
\int_S \dd\omega
=
\int_{\partial S}\omega.
\end{equation}

Формула \eqref{eq:general-stokes} утверждает: интеграл локальной производной по области равен интегралу исходной величины по границе. При \(k=1\) она превращается в формулу Ньютона--Лейбница. При \(k=2\) получаются формулы Грина и классическая формула Стокса. При \(k=3\) получается формула Гаусса--Остроградского. Различие между этими теоремами связано не с разными принципами, а только со степенью интегрируемой формы и размерностью области.

\begin{problem}[Практическая задача: проверка потенциальности поля]
В области \(D=\R^3\) задана форма работы
\[
\omega
=yz\,\dd x+xz\,\dd y+xy\,\dd z.
\]
Нужно проверить, является ли она дифференциалом некоторой функции, восстановить эту функцию и вычислить работу поля между точками
\[
A=(1,1,1),
\qquad
B=(2,3,4)
\]
без выбора конкретного пути.
\end{problem}

\begin{proof}[Решение]
Здесь
\[
P=yz,
\qquad
Q=xz,
\qquad
R=xy.
\]
По формуле \eqref{eq:exterior-derivative-one-form}
\[
\frac{\partial R}{\partial y}
-
\frac{\partial Q}{\partial z}
=x-x=0,
\]
\[
\frac{\partial P}{\partial z}
-
\frac{\partial R}{\partial x}
=y-y=0,
\]
\[
\frac{\partial Q}{\partial x}
-
\frac{\partial P}{\partial y}
=z-z=0.
\]
Следовательно, \(\dd\omega=0\). В односвязной области \(\R^3\) это условие позволяет искать потенциал \(f\), для которого \(\omega=\dd f\).

Из равенства
\[
\frac{\partial f}{\partial x}=yz
\]
получаем
\[
f(x,y,z)=xyz+C(y,z).
\]
Сравнение с условиями
\[
\frac{\partial f}{\partial y}=xz,
\qquad
\frac{\partial f}{\partial z}=xy
\]
показывает, что \(C\) является постоянной. Можно взять
\[
f(x,y,z)=xyz.
\]

По основной теореме о криволинейном интеграле работа не зависит от пути и равна разности потенциала:
\[
\int_A^B\omega
=f(B)-f(A)
=2\cdot3\cdot4-1\cdot1\cdot1
=23.
\]
Прямое интегрирование по любой кусочно гладкой кривой от \(A\) к \(B\) дало бы то же значение.
\end{proof}

% HIDDENPOWER FINAL STYLE AUDIT APPLIED
